% Source-faithful appendix generated by build_appendix.py.
% All source paths refer to members of DER_03_companion_newresults_conventions.zip.
\begingroup
\selectlanguage{english}
\providecommand{\Set}{\mathbf{Set}}
\providecommand{\Ab}{\mathbf{Ab}}
\providecommand{\Winfty}{W_\infty}
\providecommand{\Hot}{\mathbf{Hot}}
\providecommand{\Z}{\mathbb{Z}}
\providecommand{\Hom}{\operatorname{Hom}}
\providecommand{\id}{\operatorname{id}}
\providecommand{\colim}{\mathop{\operatorname{colim}}}
\newtheorem{nrdefinition}{Source definition}
\newtheorem{nrlemma}{Source lemma}
\newtheorem{nrtheorem}{Source theorem}
\newtheorem{nrproposition}{Source proposition}
\newenvironment{nrsourceclaim}[3]{\par\medskip\noindent\textbf{Source #1. #3}\quad\textit{#2}\par\smallskip}{\par\medskip}
\dgrbackchapter{Results proved along the way that we did not find in the literature}
\markright{Appendix: results and source proofs}


Some statements had to be proved along the way, while checking Grothendieck's text and the returns of the research prompts, because we could not find them in the sources we searched. They are collected here so that a reader can check them, find the literature we missed, and see where they bear on the manuscript. Our search was limited, so some of them may well be known; whoever knows a reference is asked to tell us, and the entry will cite it. No priority is claimed. Each entry gives the statement, where the proof is, who proved it (model and effort setting), whether anyone has checked it yet, and the pages of the manuscript it bears on. Statements that turned out to be known are listed at the end with the reference.

\noindent\textbf{Editorial scope.} The inventory was started on 12 November 2026 by Claude Opus~5.5 (effort max). This appendix was assembled by GPT-6 Astra (OpenAI, ChatGPT Pro mode). Each numbered section preserves the inventory statement and the available cited proof, with mathematical glyphs and scripts typeset for reading. Proofs and source comments remain attributed to their source authors. The paragraph headed ``GPT-6 Astra's reading'' is an editorial reading of the supplied argument; it is not the separate peer review requested for later work. Missing cited evidence is stated explicitly. Original source identifiers are preserved. The source extracts and the exact glyph mapping are included among the build records. No material from STRATA, PATHS, BISITE or RESIDUE is added beyond the inventory's own references.
\par\medskip
\noindent\textbf{Notation from the cited sources.} A universe $U$ is fixed; categories are small in that universe. $BX=|NX|$ denotes the classifying space, $W_\infty$ the functors inducing weak homotopy equivalences, and $\mathbf{Hot}$ the homotopy category. In $\mathrm{Cat}_{\mathrm{hot}}$, morphisms are $\pi_0\operatorname{Fun}(X,K)$. The functor $\Pi_1$ inverts every arrow, with unit $\eta_X:X\to\Pi_1X$. For a monoid $M$, $U(M)$ is its universal group. Brackets $[A,Y]$ denote free homotopy classes. Original labels such as Theorem~2.1 and G.4 refer to the cited source, not to the numbering of this appendix.


\section*{1. Finite monoids, finite dimension, and bounded homology}

\phantomsection\addcontentsline{toc}{section}{1. Finite monoids, finite dimension, and bounded homology}

\label{nr:entry1}

\noindent\textbf{Statement recorded in the inventory.}\par\smallskip

\textbf{Finite monoids: finite type, finite dimension and bounded homology coincide.} For a finite monoid M: BM has the homotopy type of a finite complex \ensuremath{\Leftrightarrow} of a finite-dimensional complex \ensuremath{\Leftrightarrow} H\textsubscript{k}(M; Z) = 0 for all large k; then the universal group U(M) is trivial. Corollaries: a "Swan theorem" for finite monoids (U(M) \ensuremath{\ne} 1 forces homology in infinitely many degrees; p-locally too); finite cohomological dimension forces U(M) = 1; with Theorem C of Ben-Bassat (arXiv:2610.38760v1) the homotopy types of finite monoids with bounded homology are exactly the finite simply connected complexes.


\noindent\textbf{Bearing on the manuscript.} V ms. 51--52 (Hot\textsubscript{0}, "types d'homotopie de type fini"); sharpens G.4 of the CATHOT return, the counterexample to Ben-Bassat's Conjecture 7.1\par\smallskip

\noindent\textbf{Proved by.} Claude Fable 5.2, effort max, 12 Nov 2026\par

\noindent\textbf{Checked by (inventory record).} not yet (Pro check prompt drafted)\par\smallskip

\noindent\textbf{Cited source.} \path{pro_reviews_max_20261112/1_CATHOT/OWN_RESULTS.md}, section 2\par\medskip

\subsection*{Statement and proof in the supplied source}

Let M be a finite monoid. Facts used, all proved in \path{G6_VERDICT.md}: \ensuremath{\pi}\textsubscript{1}(BM) = U(M); \ensuremath{\eta}: M \ensuremath{\to} U(M) is onto; for an idempotent e of the minimal ideal, H = eMe is a finite group and there is a map j\textsubscript{e}: BH \ensuremath{\to} BM (through the Karoubi envelope) inducing the surjection \ensuremath{\eta}|H on \ensuremath{\pi}\textsubscript{1}. Also H\textsubscript{k}(BM; Z) = Tor\textsubscript{k}\textsuperscript{ZM}(Z, Z) is finitely generated for every k (the bar complex is finitely generated in each degree).

\textbf{Theorem 2.1.} If H\textsubscript{k}(M; Z) = 0 for all large k, then U(M) = 1. More precisely, if a prime p divides |U(M)|, then H\textsubscript{k}(M; F\textsubscript{p}) \ensuremath{\ne} 0 for infinitely many k.

\emph{Proof.} If H\textsubscript{k}(M; Z) = 0 for k \ensuremath{\geq} n, then H\textsuperscript{q}(BM; Z) = Hom(H\textsubscript{q}, Z) \ensuremath{\oplus} Ext(H\textsubscript{q\ensuremath{-}1}, Z) = 0 for q \ensuremath{>} n. \ensuremath{\pi}\textsubscript{1}(BM) = U(M) is finite. Theorem 1.1 applied to j\textsubscript{e}: BH \ensuremath{\to} BM says \ensuremath{\eta}|H is trivial; it is onto; so U(M) = 1. For the second statement: if H\textsubscript{k}(M; F\textsubscript{p}) = 0 for all k \ensuremath{\geq} k\textsubscript{0}, then H\textsubscript{k}(M; Z) \ensuremath{\otimes} F\textsubscript{p} = 0 and Tor(H\textsubscript{k\ensuremath{-}1}(M; Z), F\textsubscript{p}) = 0 for k \ensuremath{\geq} k\textsubscript{0}, so the finitely generated groups H\textsubscript{k}(M; Z) have zero localization at p for k \ensuremath{\geq} k\textsubscript{0}, and H\textsuperscript{q}(BM; Z\textsubscript{(p)}) = 0 for q \ensuremath{>} k\textsubscript{0}. Theorem 1.2 applied to j\textsubscript{e} gives p \ensuremath{\nmid} |\ensuremath{\eta}(H)| = |U(M)|. \ensuremath{\blacksquare}

\textbf{Theorem 2.2.} For a finite monoid M the following are equivalent: (i) BM has the homotopy type of a finite CW complex; (ii) BM has the homotopy type of a finite-dimensional CW complex; (iii) H\textsubscript{k}(M; Z) = 0 for all large k. If they hold, U(M) = 1 and BM is simply connected.

\emph{Proof.} (i) \ensuremath{\Rightarrow} (ii) \ensuremath{\Rightarrow} (iii) are clear. (iii) \ensuremath{\Rightarrow} (i): by 2.1, BM is simply connected; its homology groups are finitely generated and vanish above some degree n; by Hatcher, \emph{Algebraic Topology}, Proposition 4C.1, a simply connected CW complex with finitely generated homology is homotopy equivalent to a CW complex with one k-cell for each Z summand of H\textsubscript{k} and one k-cell and one (k+1)-cell for each finite cyclic summand of H\textsubscript{k}: a finite complex of dimension \ensuremath{\leq} n + 1. \ensuremath{\blacksquare}

\textbf{Corollary 2.3 (Swan's theorem for finite monoids).} A finite monoid with U(M) \ensuremath{\ne} 1 has H\textsubscript{k}(M; Z) \ensuremath{\ne} 0 for infinitely many k; if p divides |U(M)|, then H\textsubscript{k}(M; F\textsubscript{p}) \ensuremath{\ne} 0 for infinitely many k. (For a finite group this is Swan's theorem.)

\textbf{Corollary 2.4.} If the trivial left (or right) ZM-module Z has a projective resolution of finite length (M has finite left or right cohomological dimension), then U(M) = 1 and BM has the homotopy type of a finite simply connected CW complex.

\textbf{Corollary 2.5 (with Theorem C of arXiv:2610.38760v1; recent, unrefereed).} The homotopy types of the classifying spaces of finite monoids with bounded integral homology are exactly the homotopy types of finite simply connected CW complexes.


\subsection*{GPT-6 Astra\textquotesingle s reading}

The argument explicitly connects the surjective group map $eMe\to U(M)$ to a map of classifying spaces through the Karoubi envelope; section~12 below reproduces that construction. The bounded-homology implication then uses the cohomological result reproduced in section~2, and the finite-generation input is supplied by the degreewise finite bar complex. The finite-complex conclusion invokes Hatcher's Proposition~4C.1; the positive realization assertion in Corollary~2.5 invokes Theorem~C of the stated, unrefereed preprint. The proof records these dependencies rather than proving them here. I found no additional missing inference in the displayed chain. In the universal-coefficient line the direct-sum decomposition is noncanonical, which does not affect the vanishing implication used. The computational observation following the source corollaries is not needed for the proof and is not reproduced here.

\par\bigskip

\section*{2. Maps from finite-group classifying spaces under a cohomological bound}

\phantomsection\addcontentsline{toc}{section}{2. Maps from finite-group classifying spaces under a cohomological bound}

\label{nr:entry2}

\noindent\textbf{Statement recorded in the inventory.}\par\smallskip

\textbf{Miller's Theorem B under a cohomological bound.} A map BG \ensuremath{\to} X, G finite, X connected with finite (or residually finite, or linear) \ensuremath{\pi}\textsubscript{1} and H\textsuperscript{q}(X; Z) = 0 for q \ensuremath{>} n, is trivial on \ensuremath{\pi}\textsubscript{1}; p-local form with Z\textsubscript{(p)} coefficients.


\noindent\textbf{Bearing on the manuscript.} the tool for result 1\par\smallskip

\noindent\textbf{Proved by.} Claude Fable 5.2, effort max\par

\noindent\textbf{Checked by (inventory record).} not yet\par\smallskip

\noindent\textbf{Cited source.} \path{pro_reviews_max_20261112/1_CATHOT/OWN_RESULTS.md}, section 1\par\medskip

\subsection*{Statement and proof in the supplied source}

\textbf{Theorem 1.1.} Let X be a connected CW complex with finite fundamental group \ensuremath{\pi}, and suppose H\textsuperscript{q}(X; Z) = 0 for all q \ensuremath{>} n. Let G be a finite group and f: BG \ensuremath{\to} X a map. Then f\textsubscript{*}: G \ensuremath{\to} \ensuremath{\pi} is trivial.

\textbf{Theorem 1.2 (p-local form).} Let X be a connected CW complex with finite fundamental group \ensuremath{\pi}, p a prime, and suppose H\textsuperscript{q}(X; Z\textsubscript{(p)}) = 0 for all q \ensuremath{>} n. Then for every finite group G and every map f: BG \ensuremath{\to} X, the order of the image of f\textsubscript{*}: G \ensuremath{\to} \ensuremath{\pi} is prime to p.

Miller's Theorem B (Ann. of Math. 120 (1984), p. 40) has "X finite-dimensional" as hypothesis and no condition on \ensuremath{\pi}\textsubscript{1}(X). Here the dimension hypothesis is weakened to a bound on ordinary cohomology with constant coefficients, at the price of assuming \ensuremath{\pi}\textsubscript{1}(X) finite (residually finite, or linear over C, is enough for 1.1; see the remark at the end). The proof is Miller's and Atiyah's, with the covering-space step replaced by the regular representation.

\emph{K-theory used.} K\textsuperscript{0}(Y) = [Y, Z \ensuremath{\times} BU] for CW complexes Y; it is a ring, contravariant in Y. For a CW complex Y put F\textsuperscript{q}(Y) = ker(K\textsuperscript{0}(Y) \ensuremath{\to} K\textsuperscript{0}(Y\textsuperscript{(q\ensuremath{-}1)})). (K1) F\textsuperscript{a} \ensuremath{\cdot} F\textsuperscript{b} \ensuremath{\subset} F\textsuperscript{a+b}. Proof: if x \ensuremath{\in} F\textsuperscript{a}, y \ensuremath{\in} F\textsuperscript{b}, lift them to \ensuremath{\widetilde{x}} \ensuremath{\in} K\textsuperscript{0}(Y, Y\textsuperscript{(a\ensuremath{-}1)}), \ensuremath{\widetilde y} \ensuremath{\in} K\textsuperscript{0}(Y, Y\textsuperscript{(b\ensuremath{-}1)}); the external product \ensuremath{\widetilde{x}} \ensuremath{\boxtimes} \ensuremath{\widetilde y} lies in K\textsuperscript{0}(Y \ensuremath{\times} Y, Y\textsuperscript{(a\ensuremath{-}1)} \ensuremath{\times} Y \ensuremath{\cup} Y \ensuremath{\times} Y\textsuperscript{(b\ensuremath{-}1)}), and the (a+b\ensuremath{-}1)-skeleton of Y \ensuremath{\times} Y is contained in Y\textsuperscript{(a\ensuremath{-}1)} \ensuremath{\times} Y \ensuremath{\cup} Y \ensuremath{\times} Y\textsuperscript{(b\ensuremath{-}1)}; a cellular approximation \ensuremath{\Delta}\ensuremath{\prime} of the diagonal maps Y\textsuperscript{(a+b\ensuremath{-}1)} into it, so (xy)|Y\textsuperscript{(a+b\ensuremath{-}1)} = \ensuremath{\Delta}\ensuremath{\prime}\textsuperscript{*}(x \ensuremath{\boxtimes} y)|\ldots{} = 0. (K2) If H\textsuperscript{q}(Y; Z) = 0 for all q \ensuremath{>} n, then F\textsuperscript{q}(Y) = F\textsuperscript{q+1}(Y) for all q \ensuremath{>} n; hence F\textsuperscript{n+1}(Y) = \ensuremath{\bigcap}\textsubscript{q} F\textsuperscript{q}(Y). Proof: let x \ensuremath{\in} F\textsuperscript{q}, q \ensuremath{>} n. Then x|Y\textsuperscript{(q)} comes from a class c \ensuremath{\in} K\textsuperscript{0}(Y\textsuperscript{(q)}, Y\textsuperscript{(q\ensuremath{-}1)}) = C\textsuperscript{q}(Y; \ensuremath{\widetilde{K}}\textsuperscript{0}(S\textsuperscript{q})), the cellular cochains with coefficients Z (q even) or 0 (q odd). Since x|Y\textsuperscript{(q)} extends over Y\textsuperscript{(q+1)}, the image of c in K\textsuperscript{1}(Y\textsuperscript{(q+1)}, Y\textsuperscript{(q)}) vanishes; that image is the cellular coboundary \ensuremath{\delta}c. So c is a cocycle, hence (H\textsuperscript{q} = 0) a coboundary c = \ensuremath{\delta}b with b \ensuremath{\in} K\textsuperscript{\ensuremath{-}1}(Y\textsuperscript{(q\ensuremath{-}1)}, Y\textsuperscript{(q\ensuremath{-}2)}), and the image of \ensuremath{\delta}b in K\textsuperscript{0}(Y\textsuperscript{(q)}) is zero because K\textsuperscript{\ensuremath{-}1}(Y\textsuperscript{(q\ensuremath{-}1)}) \ensuremath{\to} K\textsuperscript{0}(Y\textsuperscript{(q)}, Y\textsuperscript{(q\ensuremath{-}1)}) \ensuremath{\to} K\textsuperscript{0}(Y\textsuperscript{(q)}) is exact. So x|Y\textsuperscript{(q)} = 0, that is x \ensuremath{\in} F\textsuperscript{q+1}. (K3) A cellular map g: Y\ensuremath{\prime} \ensuremath{\to} Y satisfies g\textsuperscript{*}F\textsuperscript{q}(Y) \ensuremath{\subset} F\textsuperscript{q}(Y\ensuremath{\prime}). (K1)--(K3) hold verbatim for the p-localized theory K\textsubscript{(p)} (represented by the p-localization of the spectrum KU, coefficients Z\textsubscript{(p)} in even degrees), with H\textsuperscript{q}(\ensuremath{-}; Z\textsubscript{(p)}) in (K2); and K\textsuperscript{0}\textsubscript{(p)}(Y) = K\textsuperscript{0}(Y) \ensuremath{\otimes} Z\textsubscript{(p)} for finite Y.

\emph{Atiyah's theorem used} (Publ. Math. IHES 9 (1961)). For a finite group G let \ensuremath{\mathcal K}\textsuperscript{0}(B\textsubscript{G}) = lim\textsubscript{m} K\textsuperscript{0}(B\textsubscript{G}\textsuperscript{(m)}) over the finite skeleta of a model of BG with finite skeleta, and \ensuremath{\alpha}: R(G) \ensuremath{\to} \ensuremath{\mathcal K}\textsuperscript{0}(B\textsubscript{G}) the map sending a representation to its associated bundle. (7.2): \ensuremath{\alpha} induces an isomorphism of the I(G)-adic completion R(G)\textasciicircum{} onto \ensuremath{\mathcal K}*(B\textsubscript{G}). (6.10): the kernel of R(G) \ensuremath{\to} R(G)\textasciicircum{} is the set of virtual characters vanishing on all elements of prime-power order; (6.11): it is zero for a p-group. Two models of BG with finite skeleta are related by cellular homotopy equivalences, so "\ensuremath{\alpha}(\ensuremath{\theta}) restricts to zero on every finite skeleton" does not depend on the model.

\emph{Proof of 1.1.} Make f cellular; BG = |NG| has finite skeleta. Let \ensuremath{\rho}: \ensuremath{\pi} \ensuremath{\to} GL\textsubscript{d}(C) be the regular representation (d = |\ensuremath{\pi}|), V = \ensuremath{\widetilde{X}} \ensuremath{\times}\textsubscript{\ensuremath{\pi}} C\textsuperscript{d} the flat bundle on X, v = [V] \ensuremath{-} d \ensuremath{\in} K\textsuperscript{0}(X). Since X is connected, v \ensuremath{\in} F\textsuperscript{1}(X); by (K1), v\textsuperscript{n+1} \ensuremath{\in} F\textsuperscript{n+1}(X) = \ensuremath{\bigcap}\textsubscript{q} F\textsuperscript{q}(X) by (K2). By (K3), f\textsuperscript{*}(v\textsuperscript{n+1}) restricts to zero on every skeleton of BG. Now f\textsuperscript{*}V is the flat bundle of the representation \ensuremath{\rho}\ensuremath{\circ}\ensuremath{\varphi} of G, \ensuremath{\varphi} = f\textsubscript{*}; so \ensuremath{\alpha}((\ensuremath{\rho}\ensuremath{\varphi} \ensuremath{-} d)\textsuperscript{n+1}) = 0 in \ensuremath{\mathcal K}\textsuperscript{0}(B\textsubscript{G}). By (7.2) and (6.10) the virtual character (\ensuremath{\chi}\textsubscript{\ensuremath{\rho}\ensuremath{\varphi}} \ensuremath{-} d)\textsuperscript{n+1} vanishes at every g \ensuremath{\in} G of prime-power order: \ensuremath{\chi}\textsubscript{\ensuremath{\rho}}(\ensuremath{\varphi}(g)) = d. The regular character takes the value d only at 1, so \ensuremath{\varphi}(g) = 1 for every g of prime-power order. Every element of a finite group is a product of elements of prime-power order (its primary components). Hence \ensuremath{\varphi} = 1. \ensuremath{\blacksquare}

\emph{Proof of 1.2.} Suppose p divides |\ensuremath{\varphi}(G)|. Take u \ensuremath{\in} \ensuremath{\varphi}(G) of order p and g with \ensuremath{\varphi}(g) = u; if g has order p\textsuperscript{a} m with p \ensuremath{\nmid} m (a \ensuremath{\geq} 1) and mm\ensuremath{\prime} \ensuremath{\equiv} 1 mod p\textsuperscript{a}, then c = g\textsuperscript{mm\ensuremath{\prime}} has p-power order and \ensuremath{\varphi}(c) = u. Put C = \ensuremath{\langle}c\ensuremath{\rangle} and restrict f to BC (cellular; take for BC the infinite lens space, whose odd skeleta are lens spaces L\textsuperscript{2j+1}). As above, with K\textsubscript{(p)}: the image of v\textsuperscript{n+1} in K\textsuperscript{0}\textsubscript{(p)}(X) lies in \ensuremath{\bigcap}\textsubscript{q} F\textsuperscript{q}, so \ensuremath{\alpha}((\ensuremath{\rho}\ensuremath{\varphi}|C \ensuremath{-} d)\textsuperscript{n+1}) goes to zero in K\textsuperscript{0}(L\textsuperscript{2j+1}) \ensuremath{\otimes} Z\textsubscript{(p)} for all j. The group \ensuremath{\widetilde{K}}\textsuperscript{0}(L\textsuperscript{2j+1}) is a finite p-group (Atiyah--Hirzebruch: \ensuremath{\widetilde{H}}\textsuperscript{even}(L\textsuperscript{2j+1}; Z) consists of copies of Z/|C|), so K\textsuperscript{0}(L\textsuperscript{2j+1}) \ensuremath{\to} K\textsuperscript{0}(L\textsuperscript{2j+1}) \ensuremath{\otimes} Z\textsubscript{(p)} is injective. Hence \ensuremath{\alpha}((\ensuremath{\rho}\ensuremath{\varphi}|C \ensuremath{-} d)\textsuperscript{n+1}) = 0 in \ensuremath{\mathcal K}\textsuperscript{0}(B\textsubscript{C}), and by (7.2), (6.11) (C is a p-group; for cyclic groups Atiyah's (8.1) suffices) (\ensuremath{\rho}\ensuremath{\varphi}|C \ensuremath{-} d)\textsuperscript{n+1} = 0 in R(C). Its value at c is (\ensuremath{\chi}\textsubscript{\ensuremath{\rho}}(u) \ensuremath{-} d)\textsuperscript{n+1} = (\ensuremath{-}d)\textsuperscript{n+1} \ensuremath{\ne} 0. Contradiction. \ensuremath{\blacksquare}

\emph{Remark.} In 1.1 the regular representation can be replaced by any finite-dimensional complex representation \ensuremath{\rho} of \ensuremath{\pi}\textsubscript{1}(X), finite or not: the same proof gives \ensuremath{\rho}(f\textsubscript{*}(g)) = 1 for all g \ensuremath{\in} G (a matrix of finite order with trace equal to its size is the identity). So for arbitrary \ensuremath{\pi}\textsubscript{1}(X) the image f\textsubscript{*}(G) lies in the intersection of the kernels of all finite-dimensional complex representations, and f\textsubscript{*} is trivial if \ensuremath{\pi}\textsubscript{1}(X) is residually finite or linear.


\subsection*{GPT-6 Astra\textquotesingle s reading}

The proof states the skeletal-filtration facts it needs and derives the character obstruction from the regular representation. The essential point in (K2) is the exactness of the skeletal relative-$K$-theory sequence: the cellular coboundary representing the relative class maps to zero in absolute $K$-theory. This establishes the filtration inclusion actually used without asserting convergence of an unrestricted Atiyah--Hirzebruch spectral sequence. The $p$-local proof reduces to finite lens-space skeleta, where the reduced $K$-groups are finite $p$-groups and localization is injective. The residual-finiteness and linearity sentence uses the stated finite-dimensional representation argument. I found no gap in these deductions from the explicitly quoted Atiyah completion and character-kernel results; those external results remain cited inputs.

\par\bigskip

\section*{3. Kan--Thurston without small quotients and the ind-category statement}

\phantomsection\addcontentsline{toc}{section}{3. Kan--Thurston without small quotients and the ind-category statement}

\label{nr:entry3}

\noindent\textbf{Statement recorded in the inventory.}\par\smallskip

\textbf{A Kan--Thurston theorem without small quotients, and E.17.} Every connected CW complex W has an acyclic map BG \ensuremath{\to} W such that every homomorphism from G to a group of cardinality \ensuremath{<} \ensuremath{\lambda} factors through \ensuremath{\pi}\textsubscript{1}(W). Hence a U-small ind-object of Cat\textsubscript{hot} is W\ensuremath{\infty}-invariant iff it is ind-isomorphic to an ind-object of groupoids.


\noindent\textbf{Bearing on the manuscript.} V ms. 46 is a theorem for "local", false for "fibrant"; VII ms. 129--132 and IX ms. 22--23: invariant path ind-categories see only \ensuremath{\Pi}\textsubscript{1}\par\smallskip

\noindent\textbf{Proved by.} Claude Fable 5.2, effort max (the reduction also in the earlier review, effort medium)\par

\noindent\textbf{Checked by (inventory record).} not yet; the Kan--Thurston cardinality bound was quoted from memory\par\smallskip

\noindent\textbf{Cited source.} \path{pro_reviews_max_20261112/1_CATHOT/OWN_RESULTS.md}, section 3\par\medskip

\subsection*{Statement and proof in the supplied source}

The return leaves open (E.17): if a sequence K = (K\textsubscript{n}) in Cat\textsubscript{hot} has a W\ensuremath{\infty}-invariant presheaf H\textsubscript{K} = colim\textsubscript{n} \ensuremath{\pi}\textsubscript{0}Fun(\ensuremath{-}, K\textsubscript{n}), is K isomorphic in Ind(Cat\textsubscript{hot}) to a sequence of groupoids? It reduces this to (E.17c): H\textsubscript{K} inverts every \ensuremath{\eta}\textsubscript{X}: X \ensuremath{\to} \ensuremath{\Pi}\textsubscript{1}X. The answer is yes.

Input from the return (verified): B.7 (for every group H and every infinite cardinal \ensuremath{\lambda} \ensuremath{\geq} |H| there is a non-abelian simple, integrally acyclic group A \ensuremath{\supset} H with |A| = \ensuremath{\lambda}); A.4 and A.7; C.1, C.4; (E.17a). Input from the literature: the Kan--Thurston theorem (D. M. Kan, W. P. Thurston, \emph{Every connected space has the homology of a K(\ensuremath{\pi},1)}, Topology 15 (1976), 253--258): for every connected CW complex E there are a group T and a map BT \ensuremath{\to} E inducing isomorphisms on homology with every local coefficient system on E (an \emph{acyclic} map); T can be taken of cardinality at most max(\ensuremath{\aleph}\textsubscript{0}, number of cells of E). I could not open the paper (the publisher's open archive refused automated access); the statement is quoted from memory and agrees with the summary in Palmer--Wu, arXiv:2510.16879v1, Introduction. It is used only in the successor step for cells of dimension \ensuremath{\geq} 3 below. Standard facts on acyclic maps (Hausmann--Husemoller): the homotopy fibre of an acyclic map is acyclic; acyclic maps are stable under homotopy pull-back and composition; an acyclic map f: A \ensuremath{\to} B of connected CW complexes is onto on \ensuremath{\pi}\textsubscript{1} with perfect kernel N, and for every space Y, f\textsuperscript{*}: [B, Y] \ensuremath{\to} [A, Y] is injective with image the classes of maps g with \ensuremath{\pi}\textsubscript{1}(g)(N) = 1.

\textbf{Theorem 3.1.} Let W be a connected CW complex with one 0-cell, \ensuremath{\pi} = \ensuremath{\pi}\textsubscript{1}(W), and \ensuremath{\lambda} an infinite cardinal not less than the number of cells of W. There are a group G with |G| \ensuremath{\leq} \ensuremath{\lambda} and a map t: BG \ensuremath{\to} W such that (i) t is acyclic (an isomorphism on homology with every local coefficient system on W); in particular \ensuremath{\tau} = \ensuremath{\pi}\textsubscript{1}(t): G \ensuremath{\to} \ensuremath{\pi} is onto; (ii) every homomorphism from G to a group of cardinality \ensuremath{<} \ensuremath{\lambda} factors through \ensuremath{\tau}. Consequently, for every space Y all of whose fundamental groups have cardinality \ensuremath{<} \ensuremath{\lambda}, t\textsuperscript{*}: [W, Y] \ensuremath{\to} [BG, Y] is a bijection.

\emph{Proof.} Enumerate the cells of dimension \ensuremath{\geq} 2 by an ordinal, in non-decreasing order of dimension: W\textsubscript{1} = W\textsuperscript{(1)}, W\textsubscript{\ensuremath{\alpha}+1} = W\textsubscript{\ensuremath{\alpha}} \ensuremath{\cup}\textsubscript{a\textsubscript{\ensuremath{\alpha}}} e\textsuperscript{m\textsubscript{\ensuremath{\alpha}}+1} with a\textsubscript{\ensuremath{\alpha}}: S\textsuperscript{m\textsubscript{\ensuremath{\alpha}}} \ensuremath{\to} W\textsubscript{\ensuremath{\alpha}}, m\textsubscript{\ensuremath{\alpha}} \ensuremath{\geq} 1, unions at limit ordinals. We build groups G\textsubscript{\ensuremath{\alpha}}, injective homomorphisms G\textsubscript{\ensuremath{\alpha}} \ensuremath{\to} G\textsubscript{\ensuremath{\beta}} (\ensuremath{\alpha} \ensuremath{<} \ensuremath{\beta}) and maps t\textsubscript{\ensuremath{\alpha}}: BG\textsubscript{\ensuremath{\alpha}} \ensuremath{\to} W\textsubscript{\ensuremath{\alpha}}, with t\textsubscript{\ensuremath{\beta}}|BG\textsubscript{\ensuremath{\alpha}} = t\textsubscript{\ensuremath{\alpha}}, satisfying (i) and (ii) for (G\textsubscript{\ensuremath{\alpha}}, W\textsubscript{\ensuremath{\alpha}}), with \ensuremath{\tau}\textsubscript{\ensuremath{\alpha}} = \ensuremath{\pi}\textsubscript{1}(t\textsubscript{\ensuremath{\alpha}}): G\textsubscript{\ensuremath{\alpha}} \ensuremath{\to} \ensuremath{\pi}\textsubscript{\ensuremath{\alpha}} = \ensuremath{\pi}\textsubscript{1}(W\textsubscript{\ensuremath{\alpha}}).

\emph{Start.} G\textsubscript{1} is the free group on the 1-cells and t\textsubscript{1}: BG\textsubscript{1} \ensuremath{\to} W\textsubscript{1} a homotopy equivalence; (ii) is empty.

\emph{Successor step.} Let a = a\textsubscript{\ensuremath{\alpha}}: S\textsuperscript{m} \ensuremath{\to} W\textsubscript{\ensuremath{\alpha}}. Claim: there are a group H\ensuremath{\prime}, an acyclic map q: BH\ensuremath{\prime} \ensuremath{\to} S\textsuperscript{m} and a homomorphism u: H\ensuremath{\prime} \ensuremath{\to} G\textsubscript{\ensuremath{\alpha}} with t\textsubscript{\ensuremath{\alpha}}\ensuremath{\circ}Bu \ensuremath{\simeq} a\ensuremath{\circ}q. For m = 1: H\ensuremath{\prime} = Z, q = id, u(1) = g with \ensuremath{\tau}\textsubscript{\ensuremath{\alpha}}(g) the class of the loop a (\ensuremath{\tau}\textsubscript{\ensuremath{\alpha}} is onto). For m \ensuremath{\geq} 2: let E be the homotopy pull-back of S\textsuperscript{m} \ensuremath{\to} W\textsubscript{\ensuremath{\alpha}} \ensuremath{\leftarrow} BG\textsubscript{\ensuremath{\alpha}}; E \ensuremath{\to} S\textsuperscript{m} is acyclic and E is connected; by Kan--Thurston choose H\ensuremath{\prime} and an acyclic map BH\ensuremath{\prime} \ensuremath{\to} E; q is the composite to S\textsuperscript{m}, and the composite BH\ensuremath{\prime} \ensuremath{\to} E \ensuremath{\to} BG\textsubscript{\ensuremath{\alpha}} is homotopic to Bu for a homomorphism u. In both cases |H\ensuremath{\prime}| \ensuremath{\leq} \ensuremath{\lambda}. By B.7 choose simple acyclic groups A\ensuremath{\prime}, A\ensuremath{\prime\prime} of cardinality \ensuremath{\lambda} and embeddings i\ensuremath{\prime}: H\ensuremath{\prime} \ensuremath{\to} A\ensuremath{\prime}, i\ensuremath{\prime\prime}: H\ensuremath{\prime} \ensuremath{\to} A\ensuremath{\prime\prime}. Put P = G\textsubscript{\ensuremath{\alpha}} \ensuremath{\times} A\ensuremath{\prime} and v = (u, i\ensuremath{\prime}): H\ensuremath{\prime} \ensuremath{\to} P, which is injective, and

  G\textsubscript{\ensuremath{\alpha}+1} = P *\textsubscript{H\ensuremath{\prime}} A\ensuremath{\prime\prime}  (amalgamated along v and i\ensuremath{\prime\prime}).

P and A\ensuremath{\prime\prime} embed in the amalgam, so G\textsubscript{\ensuremath{\alpha}} \ensuremath{\to} G\textsubscript{\ensuremath{\alpha}+1}, g \ensuremath{\mapsto} (g, 1), is injective, and BG\textsubscript{\ensuremath{\alpha}+1} is the double mapping cylinder of BP \ensuremath{\leftarrow} BH\ensuremath{\prime} \ensuremath{\to} BA\ensuremath{\prime\prime}. Define t\textsubscript{\ensuremath{\alpha}+1} on it by t\textsubscript{\ensuremath{\alpha}}\ensuremath{\circ}pr\textsubscript{1} on BP, by the cone of q on the mapping cylinder of BH\ensuremath{\prime} \ensuremath{\to} BA\ensuremath{\prime\prime} (BA\ensuremath{\prime\prime} going to the centre of the new cell), and by the chosen homotopy t\textsubscript{\ensuremath{\alpha}}\ensuremath{\circ}Bu \ensuremath{\simeq} a\ensuremath{\circ}q on the other cylinder; then t\textsubscript{\ensuremath{\alpha}+1}|BG\textsubscript{\ensuremath{\alpha}} = t\textsubscript{\ensuremath{\alpha}}. \emph{Acyclicity:} for a local system L on W\textsubscript{\ensuremath{\alpha}+1} compare the Mayer--Vietoris sequences of the two double mapping cylinders. On BP: H\textsubscript{*}(G\textsubscript{\ensuremath{\alpha}} \ensuremath{\times} A\ensuremath{\prime}; L) = H\textsubscript{*}(G\textsubscript{\ensuremath{\alpha}}; L) because A\ensuremath{\prime} is acyclic and acts trivially on L (Lyndon--Hochschild--Serre), and t\textsubscript{\ensuremath{\alpha}} is acyclic. On BH\ensuremath{\prime}: q is acyclic. On BA\ensuremath{\prime\prime}: L is constant on the cell and A\ensuremath{\prime\prime} is acyclic. The five lemma gives (i). \emph{No small quotients:} let \ensuremath{\rho}: G\textsubscript{\ensuremath{\alpha}+1} \ensuremath{\to} \ensuremath{\Lambda} with |\ensuremath{\Lambda}| \ensuremath{<} \ensuremath{\lambda}. \ensuremath{\rho} is trivial on A\ensuremath{\prime\prime} and on 1 \ensuremath{\times} A\ensuremath{\prime} (simple groups of cardinality \ensuremath{\lambda}). For h \ensuremath{\in} H\ensuremath{\prime}, v(h) = (u(h), 1)(1, i\ensuremath{\prime}(h)) is identified with i\ensuremath{\prime\prime}(h) \ensuremath{\in} A\ensuremath{\prime\prime}, so \ensuremath{\rho}(u(h), 1) = 1. By (ii) for G\textsubscript{\ensuremath{\alpha}}, \ensuremath{\rho}|G\textsubscript{\ensuremath{\alpha}} = \ensuremath{\overline{\rho}}\ensuremath{\circ}\ensuremath{\tau}\textsubscript{\ensuremath{\alpha}}, and \ensuremath{\overline{\rho}} kills \ensuremath{\tau}\textsubscript{\ensuremath{\alpha}}(u(H\ensuremath{\prime})), which is trivial for m \ensuremath{\geq} 2 and is the cyclic subgroup generated by the class of the attaching loop for m = 1. By van Kampen \ensuremath{\overline{\rho}} factors through \ensuremath{\pi}\textsubscript{\ensuremath{\alpha}+1} = \ensuremath{\pi}\textsubscript{1}(W\textsubscript{\ensuremath{\alpha}+1}). Since G\textsubscript{\ensuremath{\alpha}+1} is generated by G\textsubscript{\ensuremath{\alpha}}, A\ensuremath{\prime}, A\ensuremath{\prime\prime}, and \ensuremath{\tau}\textsubscript{\ensuremath{\alpha}+1} is trivial on A\ensuremath{\prime} and A\ensuremath{\prime\prime} and equals G\textsubscript{\ensuremath{\alpha}} \ensuremath{\to} \ensuremath{\pi}\textsubscript{\ensuremath{\alpha}} \ensuremath{\to} \ensuremath{\pi}\textsubscript{\ensuremath{\alpha}+1} on G\textsubscript{\ensuremath{\alpha}}, \ensuremath{\rho} factors through \ensuremath{\tau}\textsubscript{\ensuremath{\alpha}+1}.

\emph{Limit step.} G\textsubscript{\ensuremath{\gamma}} = \ensuremath{\bigcup} G\textsubscript{\ensuremath{\alpha}}, t\textsubscript{\ensuremath{\gamma}} = \ensuremath{\bigcup} t\textsubscript{\ensuremath{\alpha}}. Homology with local coefficients commutes with these unions, and a homomorphism from G\textsubscript{\ensuremath{\gamma}} to \ensuremath{\Lambda} factors through colim \ensuremath{\pi}\textsubscript{\ensuremath{\alpha}} = \ensuremath{\pi}\textsubscript{\ensuremath{\gamma}}.

\emph{The consequence.} Let N = ker \ensuremath{\tau}. A map g: BG \ensuremath{\to} Y induces G \ensuremath{\to} \ensuremath{\pi}\textsubscript{1}(Y, g(*)), which factors through \ensuremath{\tau} by (ii); so every class in [BG, Y] kills N, and the universal property of the acyclic map t gives the bijection. \ensuremath{\blacksquare}

\textbf{Theorem 3.2 (E.17 is true).} Let I be a U-small filtered category and K: I \ensuremath{\to} Cat\textsubscript{hot} a diagram of U-small categories. Suppose H\textsubscript{K}(X) = colim\textsubscript{i} \ensuremath{\pi}\textsubscript{0}Fun(X, K\textsubscript{i}) sends every functor in W\ensuremath{\infty} between U-small categories to a bijection. Then H\textsubscript{K} sends every unit \ensuremath{\eta}\textsubscript{X}: X \ensuremath{\to} \ensuremath{\Pi}\textsubscript{1}X to a bijection; equivalently (E.17a) the canonical morphism K \ensuremath{\to} (\ensuremath{\Pi}\textsubscript{1}K\textsubscript{i})\textsubscript{i\ensuremath{\in}I} is an isomorphism in Ind(Cat\textsubscript{hot}). So a U-small ind-object of Cat\textsubscript{hot} is W\ensuremath{\infty}-invariant if and only if it is isomorphic to an ind-object of groupoids, and then H\textsubscript{K}(X) = colim\textsubscript{i} [BX, B\ensuremath{\Pi}\textsubscript{1}K\textsubscript{i}] depends only on \ensuremath{\Pi}\textsubscript{1}X.

\emph{Proof.} Let \ensuremath{\mu} be an infinite cardinal with |Mor K\textsubscript{i}| \ensuremath{\leq} \ensuremath{\mu} for all i. Let H\ensuremath{\prime}(Y) = colim\textsubscript{i} [Y, BK\textsubscript{i}] on Hot. By A.7 the canonical map c\textsubscript{X}: H\textsubscript{K}(X) \ensuremath{\to} H\ensuremath{\prime}(BX) is a bijection for every U-small X, naturally in X. Fix X. Choose a base object in each component c of X, let \ensuremath{\pi}\textsubscript{c} be the fundamental group there, and let \ensuremath{\lambda} \ensuremath{>} \ensuremath{\mu} be an infinite cardinal not less than the number of cells of BX. By 3.1 (applied to a model of the component BX\textsubscript{c} with one 0-cell) choose groups G\textsubscript{c}, surjections p\textsubscript{c}: G\textsubscript{c} \ensuremath{\to} \ensuremath{\pi}\textsubscript{c} and acyclic maps t\textsubscript{c}: BG\textsubscript{c} \ensuremath{\to} BX\textsubscript{c} inducing p\textsubscript{c}, such that every homomorphism from G\textsubscript{c} to a group of cardinality \ensuremath{<} \ensuremath{\lambda} factors through p\textsubscript{c}. Let \ensuremath{\mathcal G} = \ensuremath{\coprod}\textsubscript{c} BG\textsubscript{c} and \ensuremath{\mathcal P} = \ensuremath{\coprod}\textsubscript{c} B\ensuremath{\pi}\textsubscript{c}, with p: \ensuremath{\mathcal G} \ensuremath{\to} \ensuremath{\mathcal P} and the equivalence of groupoids \ensuremath{\varepsilon}: \ensuremath{\mathcal P} \ensuremath{\to} \ensuremath{\Pi}\textsubscript{1}X. (a) p\textsuperscript{*}: Fun(\ensuremath{\mathcal P}, K\textsubscript{i}) \ensuremath{\to} Fun(\ensuremath{\mathcal G}, K\textsubscript{i}) is an isomorphism of categories for every i: a functor \ensuremath{\mathcal G} \ensuremath{\to} K\textsubscript{i} is a family of objects k\textsubscript{c} with homomorphisms G\textsubscript{c} \ensuremath{\to} Aut(k\textsubscript{c}), and |Aut(k\textsubscript{c})| \ensuremath{\leq} \ensuremath{\mu} \ensuremath{<} \ensuremath{\lambda}, so these factor uniquely through p\textsubscript{c}; a natural transformation is a family of arrows satisfying conditions that are the same for G\textsubscript{c} and for its quotient \ensuremath{\pi}\textsubscript{c}. Hence H\textsubscript{K}(p) is a bijection, and so is H\ensuremath{\prime}(Bp) (c is natural). (b) t = \ensuremath{\coprod} t\textsubscript{c}: B\ensuremath{\mathcal G} \ensuremath{\to} BX induces bijections [BX, BK\textsubscript{i}] \ensuremath{\to} [B\ensuremath{\mathcal G}, BK\textsubscript{i}] for every i, by the last statement of 3.1 (the fundamental groups of BK\textsubscript{i} have cardinality \ensuremath{\leq} \ensuremath{\mu}). Hence H\ensuremath{\prime}(t) is a bijection. (c) B\ensuremath{\eta}\textsubscript{X}\ensuremath{\circ}t = B\ensuremath{\varepsilon}\ensuremath{\circ}Bp in Hot: both are maps from B\ensuremath{\mathcal G} into the classifying space of a groupoid inducing the same functor on fundamental groupoids up to isomorphism, and C.1 applies. So H\ensuremath{\prime}(t)\ensuremath{\circ}H\ensuremath{\prime}(B\ensuremath{\eta}\textsubscript{X}) = H\ensuremath{\prime}(Bp)\ensuremath{\circ}H\ensuremath{\prime}(B\ensuremath{\varepsilon}) is a bijection and H\ensuremath{\prime}(t) is a bijection; therefore H\ensuremath{\prime}(B\ensuremath{\eta}\textsubscript{X}) is a bijection, and by naturality of c so is H\textsubscript{K}(\ensuremath{\eta}\textsubscript{X}). The rest is (E.17a), that is, A.4 for the reflection \ensuremath{\Pi}\textsubscript{1} \ensuremath{\dashv} j of C.4. \ensuremath{\blacksquare}

\textbf{Corollary 3.3 (what this says about V ms. 44--46, VII ms. 129--132, IX ms. 22--23).} (a) Grothendieck's sentence at V ms. 46, "cela équivaut à la condition qu'il puisse se représenter par un système inductif sur N formé de catégories de Thomason--Kan", is a theorem when "catégorie de Thomason--Kan" is read as "local category" (C.3): a sequence (K\textsubscript{n}) has a W\ensuremath{\infty}-invariant functor exactly when it is ind-isomorphic to the sequence of groupoids (\ensuremath{\Pi}\textsubscript{1}K\textsubscript{n}). With "Thomason-fibrant" it is false (C.6). (b) Consequently an invariant functor only sees fundamental groupoids: H(X) = H(\ensuremath{\Pi}\textsubscript{1}X). For the path ind-categories H(a, b) = (Ch\textsubscript{n,\ensuremath{\Sigma}}(a, b))\textsubscript{n} of a U-small (M, \ensuremath{\Sigma}), the "condition préalable" of VII ms. 129--132 and test (b) of IX ms. 22, imposed on all of Cat\textsubscript{U}, forces Hom\textsubscript{M(Y)\ensuremath{\Sigma}\textsubscript{Y}\textsuperscript{-1}}(a\textsubscript{Y}, b\textsubscript{Y}) to depend only on \ensuremath{\Pi}\textsubscript{1}Y; in particular it equals Hom\textsubscript{M\ensuremath{\Sigma}\textsuperscript{-1}}(a, b) for every simply connected Y. (This uses Grothendieck's description of these Hom-sets as colim\textsubscript{n} \ensuremath{\pi}\textsubscript{0}Fun(Y, Ch\textsubscript{n,\ensuremath{\Sigma}}(a, b)), V ms. 45, which I did not re-derive.) So for small (M, \ensuremath{\Sigma}) the condition can hold only when the derived mapping spaces are, as ind-objects, 1-types: a precise form of his remark that the condition "n'est pratiquement jamais satisfaite" (VII ms. 132) and of his doubt at IX ms. 23. (c) The filtered case of D.8 of the return (the functor [B\ensuremath{-}, Y] is a U-small filtered colimit of representables only if Y is a 1-type) follows: if [B\ensuremath{-}, Y] = H\textsubscript{K}, then by 3.2 the map [P\textsubscript{1}Y\ensuremath{\prime}, Y] \ensuremath{\to} [Y\ensuremath{\prime}, Y] is bijective for all Y\ensuremath{\prime}; for Y\ensuremath{\prime} = Y the identity of Y factors through Y \ensuremath{\to} P\textsubscript{1}Y, so Y is a retract in Hot of a 1-type, hence a 1-type. (d) The size restriction on I is necessary (D.7 of the return).


\subsection*{GPT-6 Astra\textquotesingle s reading}

The source supplies the amalgamated-product construction and the factorization of every homomorphism to a smaller group. It explicitly identifies two places needing a careful independent check: the successor-step comparison with local coefficients, and the cardinal bound in the Kan--Thurston input. The cited cardinal estimate was recalled from memory, and its proof is not present in the supplied files. The construction also treats compatible classifying-space models and limit unions without spelling out their replacement maps. Thus the written argument has a specific citation and model-compatibility gap at the claimed cardinal bound and its repeated use. The deduction of Theorem~3.2 from Theorem~3.1 is written out, including the functor-category isomorphism induced by $G_c\twoheadrightarrow\pi_c$, but it also cites the return's localization results A.4, A.7 and E.17a. Corollary~3.3(b) expressly relies on the manuscript's path-category formula without re-deriving it. These are the source's dependencies; this reading does not supply missing proofs or promote its conclusion to an independently checked theorem.

\par\bigskip

\section*{4. Uncountably many homotopy classes of maps to the two-sphere}

\phantomsection\addcontentsline{toc}{section}{4. Uncountably many homotopy classes of maps to the two-sphere}

\label{nr:entry4}

\noindent\textbf{Statement recorded in the inventory.}\par\smallskip

\textbf{Uncountably many maps to S\textsuperscript{2}.} If X is connected with H\textsuperscript{1}(X; Z) = 0 and S\textsuperscript{n} is a homotopy retract of X for infinitely many n with \ensuremath{\pi}\textsubscript{n}(S\textsuperscript{3}) \ensuremath{\ne} 0, then [X, S\textsuperscript{2}] is uncountable.


\noindent\textbf{Bearing on the manuscript.} D.6 of CATHOT; V ms. 43--46 (with hunch H3)\par\smallskip

\noindent\textbf{Proved by.} Claude Fable 5.2, effort max\par

\noindent\textbf{Checked by (inventory record).} not yet\par\smallskip

\noindent\textbf{Cited source.} \path{pro_reviews_max_20261112/1_CATHOT/OWN_RESULTS.md}, section 4\par\medskip

\subsection*{Statement and proof in the supplied source}

\textbf{Proposition 4.1.} Let X be a connected CW complex with H\textsuperscript{1}(X; Z) = 0. Suppose that for infinitely many n with \ensuremath{\pi}\textsubscript{n}(S\textsuperscript{3}) \ensuremath{\ne} 0 the sphere S\textsuperscript{n} is a homotopy retract of X: there are maps s\textsubscript{n}: S\textsuperscript{n} \ensuremath{\to} X and f\textsubscript{n}: X \ensuremath{\to} S\textsuperscript{n} with f\textsubscript{n}\ensuremath{\circ}s\textsubscript{n} of degree \ensuremath{\pm}1. Then [X, S\textsuperscript{3}] and [X, S\textsuperscript{2}] are uncountable.

(\ensuremath{\pi}\textsubscript{n}(S\textsuperscript{3}) \ensuremath{\ne} 0 for n = 2p, p prime, by Serre: the p-primary component of \ensuremath{\pi}\textsubscript{2p}(S\textsuperscript{3}) is Z/p.)

\emph{Proof.} Let N be the set of such n. Take f\textsubscript{n} and s\textsubscript{n} cellular for the cell structure of S\textsuperscript{n} with one 0-cell and one n-cell; then f\textsubscript{n}(X\textsuperscript{(n\ensuremath{-}1)}) is the base point. Choose pointed maps a\textsubscript{n}: S\textsuperscript{n} \ensuremath{\to} S\textsuperscript{3} representing non-zero classes, with base point 1 \ensuremath{\in} S\textsuperscript{3} = SU(2), and put g\textsubscript{n} = a\textsubscript{n}\ensuremath{\circ}f\textsubscript{n}. For A \ensuremath{\subset} N define g\textsubscript{A}(x) = \ensuremath{\prod}\textsubscript{n\ensuremath{\in}A} g\textsubscript{n}(x), the product in increasing order in the group S\textsuperscript{3}; on X\textsuperscript{(k)} only the factors with n \ensuremath{\leq} k differ from 1, so g\textsubscript{A} is continuous. Let A \ensuremath{\ne} A\ensuremath{\prime} and n = min(A \ensuremath{\triangle} A\ensuremath{\prime}), say n \ensuremath{\in} A. On X\textsuperscript{(n)}: g\textsubscript{A} = h\ensuremath{\cdot}g\textsubscript{n} and g\textsubscript{A\ensuremath{\prime}} = h, with h the product of the common factors of index \ensuremath{<} n. If g\textsubscript{A} \ensuremath{\simeq} g\textsubscript{A\ensuremath{\prime}}, then h\ensuremath{\cdot}g\textsubscript{n} \ensuremath{\simeq} h on X\textsuperscript{(n)}, hence g\textsubscript{n}|X\textsuperscript{(n)} is null-homotopic (multiply pointwise by h\textsuperscript{-1}), hence g\textsubscript{n}\ensuremath{\circ}s\textsubscript{n} = a\textsubscript{n}\ensuremath{\circ}f\textsubscript{n}\ensuremath{\circ}s\textsubscript{n} is null-homotopic; but it represents \ensuremath{\pm} the class of a\textsubscript{n} (free and pointed classes agree in S\textsuperscript{3}). So A \ensuremath{\mapsto} [g\textsubscript{A}] is injective on the subsets of N. For S\textsuperscript{2}: the Hopf map S\textsuperscript{3} \ensuremath{\to} S\textsuperscript{2} is a principal S\textsuperscript{1}-bundle, so two maps X \ensuremath{\to} S\textsuperscript{3} with homotopic images in S\textsuperscript{2} differ, after a homotopy, by a map X \ensuremath{\to} S\textsuperscript{1}, which is null-homotopic because H\textsuperscript{1}(X; Z) = 0; thus [X, S\textsuperscript{3}] \ensuremath{\to} [X, S\textsuperscript{2}] is injective. \ensuremath{\blacksquare}


\subsection*{GPT-6 Astra\textquotesingle s reading}

The infinite product is locally finite on every CW skeleton, so the weak topology of the CW complex supplies its continuity. The least differing index compares the two maps on the corresponding skeleton; the cellular representative of $s_n$ lands there, and multiplication by the inverse of the common initial product gives the stated contradiction. The Hopf-bundle lifting argument then identifies two lifts differing by a map to $S^1$, whose class vanishes by $H^1(X;\mathbb Z)=0$. I found the proof complete from the stated hypotheses and standard cellular approximation and bundle homotopy lifting. It does not construct a finite category satisfying the retract hypothesis; the unlisted speculative finite-category consequence is not added to this appendix.

\par\bigskip

\section*{5. The thirty-four-arrow lower bound for the projective plane}

\phantomsection\addcontentsline{toc}{section}{5. The thirty-four-arrow lower bound for the projective plane}

\label{nr:entry5}

\noindent\textbf{Statement recorded in the inventory.}\par\smallskip

\textbf{Posets of type RP\textsuperscript{2} need at least 34 arrows} (the return proved 28).


\noindent\textbf{Bearing on the manuscript.} CATHOT G.8\par\smallskip

\noindent\textbf{Proved by.} Claude Fable 5.2, effort max\par

\noindent\textbf{Checked by (inventory record).} not yet\par\smallskip

\noindent\textbf{Cited source.} \path{pro_reviews_max_20261112/1_CATHOT/OWN_RESULTS.md}, section 5\par\medskip

\subsection*{Statement and proof in the supplied source}

The return (G.8) proves 28 \ensuremath{\leq} a\textsubscript{pos}(RP\textsuperscript{2}) \ensuremath{\leq} 49 for the least number of arrows (pairs x \ensuremath{\leq} y) of a poset whose order complex has the homotopy type of RP\textsuperscript{2}.

\textbf{Proposition 5.1.} a\textsubscript{pos}(RP\textsuperscript{2}) \ensuremath{\geq} 34.

\emph{Proof.} Let P be such a poset with the least number of arrows; it is finite. As in G.8, its comparability graph has no vertex of degree \ensuremath{\leq} 1 (a vertex with one neighbour is a beat point and can be removed), and P has at least 13 points (Cianci--Ottina, arXiv:1512.06088v1, Corollary 5.8). If P has exactly 13 points it has 49 arrows (their Theorem 5.10). Otherwise f\textsubscript{0} \ensuremath{\geq} 14. The order complex \ensuremath{\Delta} is a flag complex with \ensuremath{\widetilde{H}}\textsubscript{2}(\ensuremath{\Delta}; F\textsubscript{2}) \ensuremath{\ne} 0. By the argument of G.7 over F\textsubscript{2} (take a full subcomplex L minimal with \ensuremath{\widetilde{H}}\textsubscript{2}(L; F\textsubscript{2}) \ensuremath{\ne} 0; Mayer--Vietoris gives \ensuremath{\widetilde{H}}\textsubscript{1}(lk\textsubscript{L} v; F\textsubscript{2}) \ensuremath{\ne} 0 for each vertex v of L, and links in flag complexes are flag), every vertex of L has at least 4 neighbours in L and L has at least 6 vertices. Summing degrees: 2f\textsubscript{1} \ensuremath{\geq} 4|L| + 2(f\textsubscript{0} \ensuremath{-} |L|), so f\textsubscript{1} \ensuremath{\geq} f\textsubscript{0} + |L| \ensuremath{\geq} f\textsubscript{0} + 6, and the number of arrows is f\textsubscript{0} + f\textsubscript{1} \ensuremath{\geq} 2f\textsubscript{0} + 6 \ensuremath{\geq} 34. \ensuremath{\blacksquare}


\subsection*{GPT-6 Astra\textquotesingle s reading}

The bound combines the cited thirteen-point classification with a degree count that is explicitly proved. Minimality in arrow count removes degree-one beat points. A vertex-minimal full subcomplex with nonzero second mod-two homology has at least six vertices and degree at least four at each of them; the remaining comparability vertices have degree at least two. The sum gives $f_0+f_1\ge 2f_0+6\ge34$ outside the thirteen-point case. The argument is complete once the cited Cianci--Ottina object bound and classification are supplied. It proves the lower bound recorded here and does not identify the global arrow minimum.

\par\bigskip

\section*{6. The reported census of small categories}

\phantomsection\addcontentsline{toc}{section}{6. The reported census of small categories}

\label{nr:entry6}

\noindent\textbf{Statement recorded in the inventory.}\par\smallskip

\textbf{Census of small categories.} The fewest arrows for finite categories of types S\textsuperscript{1}, S\textsuperscript{2}, S\textsuperscript{3} are 4, 5, 7; S\textsuperscript{1}: only the parallel pair; S\textsuperscript{2}: only the monoid P\textsubscript{2}; every smallest S\textsuperscript{3} is a monoid. Grothendieck's own S\textsuperscript{2} model (XII app. ms. 10--11) has 10 arrows.


\noindent\textbf{Bearing on the manuscript.} XII app. ms. 8--11\par\smallskip

\noindent\textbf{Proved by.} Claude Fable 5.2, effort max (exhaustive enumeration)\par

\noindent\textbf{Checked by (inventory record).} not yet\par\smallskip

\noindent\textbf{Cited source.} \path{pro_reviews_max_20261112/1_CATHOT/OWN_RESULTS.md}, section 6\par\medskip

\subsection*{Statement and proof in the supplied source}

\textbf{Statement 6.1.} (a) The numbers of monoid structures on \{0, \ldots{}, n\ensuremath{-}1\} with identity 0 are 1, 2, 11, 156, 4,122, 208,672 for n = 1, \ldots{}, 6, in 1, 2, 6, 27, 156, 1,373 classes up to isomorphism and anti-isomorphism. (b) The triples (H\textsubscript{1}, H\textsubscript{2}, H\textsubscript{3}) of BM that occur for monoids of order 6 are (0,0,0) [1,159 classes], (0,Z,0) [2], (0,Z,Z) [1], (0,Z,Z/2) [2], (0,Z/2,0) [1], (Z/2,0,Z/2) [172], (Z/3,0,Z/3) [22], (Z/4,0,Z/4) [6], (Z/5,0,Z/5) [1], (Z/6,0,Z/6) [1], (Z/2,0,Z/6) [1], ((Z/2)\textsuperscript{2},Z/2,(Z/2)\textsuperscript{3}) [5]; for order 5: (0,0,0) [126], (0,Z,0) [1, the monoid P\textsubscript{2}], (Z/2,0,Z/2) [22], (Z/3,0,Z/3) [4], (Z/4,0,Z/4) [1], (Z/5,0,Z/5) [1], ((Z/2)\textsuperscript{2},Z/2,(Z/2)\textsuperscript{3}) [1]. (c) Connected categories with k \ensuremath{\geq} 2 objects and N arrows, as labelled tables with arrows sorted by (source, target): N = 3: 2; N = 4: 11; N = 5: 76 + 6 (k = 2, 3); N = 6: 977 + 54; N = 7: 22,798 + 498 + 32 (k = 2, 3, 4). None has (H\textsubscript{1}, H\textsubscript{2}, H\textsubscript{3}) = (0, 0, Z); none with N \ensuremath{\leq} 5 has H\textsubscript{2} = Z; exactly the two labellings of the parallel pair have (Z, 0, 0) with N = 4; none with N \ensuremath{\leq} 7 has (Z/2, 0, 0). (d) Hence: the categories of type S\textsuperscript{1}, S\textsuperscript{2}, S\textsuperscript{3} with the fewest arrows have 4, 5, 7 arrows; for S\textsuperscript{1} it is the parallel pair, for S\textsuperscript{2} the monoid P\textsubscript{2} (unique up to isomorphism), and for S\textsuperscript{3} every such category is a monoid of order 7 (P\textsubscript{3} is one). (e) The six classes of monoids of order 6 with H\textsubscript{1} = 0 and H\textsubscript{2} \ensuremath{\ne} 0 have mod-p Betti numbers in degrees 1--5 as listed in \path{work/code/higher_homology_order6.out}; none has the Betti numbers of CP\textsuperscript{\ensuremath{\infty}} (0,1,0,1,0).

Method: backtracking over composition tables with associativity checks (\path{catlib.py}), canonical forms under relabelling and transposition, exact Smith normal forms of the normalized nerve boundaries (python-flint). Checks of the method: the counts in (a) are the known ones; group homology of C\textsubscript{k}, C\textsubscript{2}\textsuperscript{2}, S\textsubscript{3}; the nerves of the parallel pair, of Grothendieck's category Z of XII app. ms. 10--11 (homology of S\textsuperscript{2}), of R\textsubscript{2} (RP\textsuperscript{2}), of P\textsubscript{2} and P\textsubscript{3}.


\subsection*{GPT-6 Astra\textquotesingle s reading}

The supplied file contains the census counts and the algorithm description, but the referenced programs, enumeration certificates and output files are absent from the archive. Consequently the exhaustive lower bounds and uniqueness claims cannot be reproduced or independently certified from these attachments. The homology tests described are sufficient to exclude the listed sphere types when the enumeration is exhaustive; homology alone would not certify the positive homotopy types. The source points to separately proved models $P_2$, $P_3$ and the parallel pair for existence, but does not include an enumeration certificate in this section. The asserted ten-arrow size of Grothendieck's model is recorded by NEW\_RESULTS; this section's source statement does not give its complete composition table. This entry is therefore an exact account of a reported computational result with a missing computational proof record, not a fresh exhaustive verification.

\par\bigskip

\section*{7. The reversible localizer defined by two-local spheres}

\phantomsection\addcontentsline{toc}{section}{7. The reversible localizer defined by two-local spheres}

\label{nr:entry7}

\noindent\textbf{Statement recorded in the inventory.}\par\smallskip

\textbf{A reversible fundamental localizer not detected by spectra.} W\textsubscript{S} (homotopy cofibre null for all 2-local spheres S\textsuperscript{n}, n \ensuremath{\geq} 2) satisfies W(7), contains RP\textsuperscript{\ensuremath{\infty}} \ensuremath{\to} e and not RP\textsuperscript{2} \ensuremath{\to} e, while every reversible localizer containing RP\textsuperscript{\ensuremath{\infty}} \ensuremath{\to} e has RP\textsuperscript{2} in its stable closure.


\noindent\textbf{Bearing on the manuscript.} POLES F.14 (1); XII §5 and XVI §3: the reversibility axiom does not single out homology theories\par\smallskip

\noindent\textbf{Proved by.} Claude Fable 5.2, effort max\par

\noindent\textbf{Checked by (inventory record).} not yet (prompt REVERSIBLE\textsubscript{LOCALIZERS} drafted)\par\smallskip

\noindent\textbf{Cited source.} \path{pro_reviews_max_20261112/DERIVATEURS_EXPERT/RESULT_REVERSIBLE_NOT_STABLE.md}, sections 1--4\par\medskip

\subsection*{Statement and proof in the supplied source}

\subsection*{1. Statement}

Notation as in POLES: a fundamental localizer is a class W of functors satisfying identities, two out of three, the retraction rule, the terminal-object axiom and the relative Theorem A. W is \emph{reversible} when it satisfies Loc (6\ensuremath{\prime} b) (XVI ms. 17\ensuremath{\prime}; W(7 bis) "\ensuremath{\Leftarrow}" of XII ms. 76), equivalently (POLES B.5) when g \ensuremath{\in} W \ensuremath{\Leftrightarrow} C(g) is W-aspherical, C(g) the unreduced homotopy cofibre of Bg, pointed at the cone point. For a localizing subcategory L of Ho(Sp), R(L) = \{g : \ensuremath{\Sigma}\textsuperscript{\ensuremath{\infty}} C(g) \ensuremath{\in} L\} (POLES F.2), and T(W) = loc\{\ensuremath{\Sigma}\textsuperscript{\ensuremath{\infty}} C(g) : g \ensuremath{\in} W\}.

Let S\textsubscript{n} denote the 2-localization of the sphere S\textsuperscript{n}, and let

\begin{quote}
\ensuremath{\mathfrak N} = \{ pointed spaces X : Map\textsubscript{*}(X, S\textsubscript{n}) is weakly contractible for every n \ensuremath{\geq} 2 \},
W\textsubscript{S} = \{ functors g : C(g) \ensuremath{\in} \ensuremath{\mathfrak N} \}.
\end{quote}

\textbf{Theorem.} (a) W\textsubscript{S} is a fundamental localizer, it is reversible (so it satisfies Loc (6 a, b, c), W(7), W(7 bis)), and it is not Fl(Cat). (b) RP\textsuperscript{\ensuremath{\infty}} \ensuremath{\to} e is in W\textsubscript{S} and RP\textsuperscript{2} \ensuremath{\to} e is not. (c) Every reversible fundamental localizer W that contains RP\textsuperscript{\ensuremath{\infty}} \ensuremath{\to} e has \ensuremath{\Sigma}\textsuperscript{\ensuremath{\infty}} RP\textsuperscript{2} \ensuremath{\in} T(W), so RP\textsuperscript{2} \ensuremath{\to} e \ensuremath{\in} RT(W). Hence W\textsubscript{S} \ensuremath{\ne} RT(W\textsubscript{S}): W\textsubscript{S} is not of the form R(L), and in particular is not a homology localizer W\textsubscript{E}. The same holds for every reversible localizer between the one generated by RP\textsuperscript{\ensuremath{\infty}} \ensuremath{\to} e and W\textsubscript{S}.

So the reversible cofibration axiom of Grothendieck does not force a localizer to be "stable" or homological.

\subsection*{2. The class \ensuremath{\mathfrak N} is closed under desuspension}

\textbf{Lemma 1.} If \ensuremath{\Sigma}X \ensuremath{\in} \ensuremath{\mathfrak N} then X \ensuremath{\in} \ensuremath{\mathfrak N}.

\emph{Proof.} James: for every n \ensuremath{\geq} 1 there is, after localization at 2, a fibre sequence S\textsuperscript{n} \ensuremath{\to} \ensuremath{\Omega}S\textsuperscript{n+1} \ensuremath{\to} \ensuremath{\Omega}S\textsuperscript{2n+1} (E followed by the Hopf--James invariant H); so for n \ensuremath{\geq} 2 a fibre sequence

\begin{quote}
\ensuremath{\Omega}\textsuperscript{2} S\textsubscript{2n+1} --P--\ensuremath{>} S\textsubscript{n} --E--\ensuremath{>} \ensuremath{\Omega} S\textsubscript{n+1}.
\end{quote}

Let \ensuremath{\varphi}: X \ensuremath{\to} S\textsubscript{n} be pointed. E\ensuremath{\varphi} is adjoint to a map \ensuremath{\Sigma}X \ensuremath{\to} S\textsubscript{n+1}, which is null because Map\textsubscript{*}(\ensuremath{\Sigma}X, S\textsubscript{n+1}) is weakly contractible. So \ensuremath{\varphi} \ensuremath{\simeq} P\ensuremath{\psi} with \ensuremath{\psi}: X \ensuremath{\to} \ensuremath{\Omega}\textsuperscript{2} S\textsubscript{2n+1}, and \ensuremath{\psi} is adjoint to an element of \ensuremath{\pi}\textsubscript{1}(Map\textsubscript{*}(\ensuremath{\Sigma}X, S\textsubscript{2n+1}), const) = 0. So \ensuremath{\varphi} is null: \ensuremath{\pi}\textsubscript{0} Map\textsubscript{*}(X, S\textsubscript{n}) = *. For j \ensuremath{\geq} 1, \ensuremath{\pi}\textsubscript{j}(Map\textsubscript{*}(X, S\textsubscript{n}), const) = \ensuremath{\pi}\textsubscript{j\ensuremath{-}1}(Map\textsubscript{*}(\ensuremath{\Sigma}X, S\textsubscript{n}), const) = 0. \ensuremath{\blacksquare}

\textbf{Lemma 2.} (i) \ensuremath{\mathfrak N} is closed under pointed homotopy colimits and retracts and contains the contractible spaces. (ii) For a cofibre sequence X \ensuremath{\to} Y \ensuremath{\to} C of pointed spaces, if two of X, Y, C are in \ensuremath{\mathfrak N}, so is the third. (iii) X \ensuremath{\in} \ensuremath{\mathfrak N} \ensuremath{\Leftrightarrow} \ensuremath{\Sigma}X \ensuremath{\in} \ensuremath{\mathfrak N}. (iv) If X \ensuremath{\in} \ensuremath{\mathfrak N} then X \ensuremath{\wedge} Y\_+ \ensuremath{\in} \ensuremath{\mathfrak N} and, if also Y \ensuremath{\in} \ensuremath{\mathfrak N}, X \ensuremath{\times} Y \ensuremath{\in} \ensuremath{\mathfrak N}.

\emph{Proof.} (i) Map\textsubscript{*}(hocolim X\textsubscript{i}, Z) = holim Map\textsubscript{*}(X\textsubscript{i}, Z). (ii) Use the fibre sequences Map\textsubscript{*}(C, Z) \ensuremath{\to} Map\textsubscript{*}(Y, Z) \ensuremath{\to} Map\textsubscript{*}(X, Z) and Map\textsubscript{*}(\ensuremath{\Sigma}X, Z) \ensuremath{\to} Map\textsubscript{*}(C, Z) \ensuremath{\to} Map\textsubscript{*}(Y, Z), with Map\textsubscript{*}(\ensuremath{\Sigma}X, Z) = \ensuremath{\Omega} Map\textsubscript{*}(X, Z): X, C \ensuremath{\in} \ensuremath{\mathfrak N} \ensuremath{\Rightarrow} Y \ensuremath{\in} \ensuremath{\mathfrak N}; X, Y \ensuremath{\in} \ensuremath{\mathfrak N} \ensuremath{\Rightarrow} C \ensuremath{\in} \ensuremath{\mathfrak N}; Y, C \ensuremath{\in} \ensuremath{\mathfrak N} \ensuremath{\Rightarrow} \ensuremath{\Sigma}X \ensuremath{\in} \ensuremath{\mathfrak N} (from C \ensuremath{\to} \ensuremath{\Sigma}X \ensuremath{\to} \ensuremath{\Sigma}Y), and then X \ensuremath{\in} \ensuremath{\mathfrak N} by Lemma 1. (iii) \ensuremath{\Sigma}X is the cofibre of X \ensuremath{\to} *; Lemma 1. (iv) Map\textsubscript{*}(X \ensuremath{\wedge} Y\_+, Z) = Map(Y, Map\textsubscript{*}(X, Z)); then Y \ensuremath{\to} X \ensuremath{\times} Y \ensuremath{\to} X \ensuremath{\wedge} Y\_+. \ensuremath{\blacksquare}

\subsection*{3. W\textsubscript{S} is a reversible fundamental localizer}

For f: X \ensuremath{\to} Y, g: Y \ensuremath{\to} Z there is a cofibre sequence C(f) \ensuremath{\to} C(gf) \ensuremath{\to} C(g) whose connecting map is C(g) \ensuremath{\to} \ensuremath{\Sigma}Y \ensuremath{\to} \ensuremath{\Sigma}C(f).

\begin{itemize}
\item Identities: C(id) is a cone. Two out of three: Lemma 2 (ii).
\item Retraction rule. Let ri = id\textsubscript{X}, ir \ensuremath{\in} W\textsubscript{S}, D = C(i). From C(i) \ensuremath{\to} C(ri) \ensuremath{\simeq} * \ensuremath{\to} C(r): C(r) \ensuremath{\simeq} \ensuremath{\Sigma}D. In
  C(r) \ensuremath{\to} C(ir) \ensuremath{\to} C(i) --\ensuremath{\delta}--\ensuremath{>} \ensuremath{\Sigma}C(r) the map \ensuremath{\delta} factors through the connecting map \ensuremath{\partial}: C(i) \ensuremath{\to} \ensuremath{\Sigma}X of i, and
  \ensuremath{\partial} \ensuremath{\simeq} (\ensuremath{\Sigma}r)(\ensuremath{\Sigma}i)\ensuremath{\partial} \ensuremath{\simeq} * because (\ensuremath{\Sigma}i)\ensuremath{\partial} \ensuremath{\simeq} *. So \ensuremath{\Sigma}C(ir) \ensuremath{\simeq} \ensuremath{\Sigma}\textsuperscript{2}D \ensuremath{\vee} \ensuremath{\Sigma}D. As C(ir) \ensuremath{\in} \ensuremath{\mathfrak N}, the retract \ensuremath{\Sigma}D is in \ensuremath{\mathfrak N},
  so D \ensuremath{\in} \ensuremath{\mathfrak N} (Lemma 1): i \ensuremath{\in} W\textsubscript{S}, and r \ensuremath{\in} W\textsubscript{S} by two out of three.
\item A category with a terminal object has contractible nerve.
\item Relative Theorem A. For u: A \ensuremath{\to} B over C with all u/c \ensuremath{\in} W\textsubscript{S}: BA \ensuremath{\simeq} hocolim\textsubscript{c} B(A/c),
  BB \ensuremath{\simeq} hocolim\textsubscript{c} B(B/c) naturally, and C(Bu) is the pointed homotopy colimit over C of the C(B(u/c))
  (the unpointed one, with the section BC of cone points collapsed). Lemma 2 (i).
\item Reversible: X \ensuremath{\to} e \ensuremath{\in} W\textsubscript{S} \ensuremath{\Leftrightarrow} \ensuremath{\Sigma}X \ensuremath{\in} \ensuremath{\mathfrak N} \ensuremath{\Leftrightarrow} X \ensuremath{\in} \ensuremath{\mathfrak N} (Lemma 2 (iii)); so the W\textsubscript{S}-aspherical categories are those
  with nerve in \ensuremath{\mathfrak N} (for any base point), and g \ensuremath{\in} W\textsubscript{S} \ensuremath{\Leftrightarrow} C(g) \ensuremath{\in} \ensuremath{\mathfrak N} \ensuremath{\Leftrightarrow} C(g) is W\textsubscript{S}-aspherical. By POLES B.5
  this is Loc (6 b), and Loc (6 c), W(7) follow.
\item S\textsuperscript{2} \ensuremath{\to} e \ensuremath{\notin} W\textsubscript{S} (the identity of S\textsubscript{2}), so W\textsubscript{S} \ensuremath{\ne} Fl(Cat). Acyclic spaces are in \ensuremath{\mathfrak N} (the S\textsubscript{n} are simply
  connected, hence HZ-local), in accordance with W\textsubscript{Z} \ensuremath{\subset} W\textsubscript{S} (POLES B.6).
\end{itemize}

\subsection*{4. RP\textsuperscript{\ensuremath{\infty}}, RP\textsuperscript{2}, and the stable side}

(b) Miller's theorem (the Sullivan conjecture): Map\textsubscript{*}(BG, K) is weakly contractible for G finite and K a finite complex; for K = S\textsuperscript{n} and the 2-localization S\textsubscript{n} use the arithmetic square (RP\textsuperscript{\ensuremath{\infty}} is rationally acyclic and maps from BZ/2 into a simply connected space see only its 2-completion). So RP\textsuperscript{\ensuremath{\infty}} \ensuremath{\in} \ensuremath{\mathfrak N}. On the other hand [RP\textsuperscript{2}, S\textsubscript{2}] = H\textsuperscript{2}(RP\textsuperscript{2}; Z\textsubscript{(2)}) = Z/2 \ensuremath{\ne} 0, so RP\textsuperscript{2} \ensuremath{\notin} \ensuremath{\mathfrak N}.

(c) Let W be reversible with RP\textsuperscript{\ensuremath{\infty}} \ensuremath{\to} e \ensuremath{\in} W, A its class of aspherical pointed spaces, L = T(W). Then \ensuremath{\Sigma}\textsuperscript{\ensuremath{\infty}} A \ensuremath{\subset} L, L is a tensor ideal (POLES F.13), and: 1. \ensuremath{\Sigma}\textsuperscript{\ensuremath{\infty}} K(Z/2, n) \ensuremath{\in} L for all n: K(Z/2, n+1) is the bar construction of K(Z/2, n), a pointed homotopy    colimit of finite products, and \ensuremath{\Sigma}\textsuperscript{\ensuremath{\infty}}(X \ensuremath{\times} Y) = \ensuremath{\Sigma}\textsuperscript{\ensuremath{\infty}}X \ensuremath{\vee} \ensuremath{\Sigma}\textsuperscript{\ensuremath{\infty}}Y \ensuremath{\vee} \ensuremath{\Sigma}\textsuperscript{\ensuremath{\infty}}(X \ensuremath{\wedge} Y). Hence    HZ/2 = colim\textsubscript{n} \ensuremath{\Sigma}\textsuperscript{\ensuremath{-}n} \ensuremath{\Sigma}\textsuperscript{\ensuremath{\infty}} K(Z/2, n) \ensuremath{\in} L. 2. Snaith: \ensuremath{\Sigma}\textsuperscript{\ensuremath{\infty}} Q(RP\textsuperscript{\ensuremath{\infty}}) \ensuremath{\simeq} \ensuremath{\bigvee}\textsubscript{j} \ensuremath{\Sigma}\textsuperscript{\ensuremath{\infty}} D\textsubscript{j}(RP\textsuperscript{\ensuremath{\infty}}), and \ensuremath{\Sigma}\textsuperscript{\ensuremath{\infty}} D\textsubscript{j} X = (\ensuremath{\Sigma}\textsuperscript{\ensuremath{\infty}} X)\textsuperscript{\ensuremath{\wedge} j}\textsubscript{h\ensuremath{\Sigma}\textsubscript{j}} \ensuremath{\in} L. Kahn--Priddy: the    2-localization of Q\textsubscript{0} S\textsuperscript{0} is a retract, as a space, of Q(RP\textsuperscript{\ensuremath{\infty}}). So \ensuremath{\Sigma}\textsuperscript{\ensuremath{\infty}} (Q\textsubscript{0} S\textsuperscript{0})\textsubscript{(2)} \ensuremath{\in} L. Its iterated    bar constructions are the spaces of the spectrum \ensuremath{\tau}\textsubscript{\ensuremath{\geq}1} S\textsubscript{(2)}; as in 1, \ensuremath{\tau}\textsubscript{\ensuremath{\geq}1} S\textsubscript{(2)} \ensuremath{\in} L. 3. Smashing \ensuremath{\tau}\textsubscript{\ensuremath{\geq}1} S\textsubscript{(2)} \ensuremath{\to} S\textsubscript{(2)} \ensuremath{\to} HZ\textsubscript{(2)} with the Moore spectrum: (\ensuremath{\tau}\textsubscript{\ensuremath{\geq}1} S\textsubscript{(2)})/2 \ensuremath{\to} S/2 \ensuremath{\to} HZ/2. So    S/2 \ensuremath{\in} L, i.e. \ensuremath{\Sigma}\textsuperscript{\ensuremath{\infty}} RP\textsuperscript{2} \ensuremath{\in} L, and RP\textsuperscript{2} \ensuremath{\to} e \ensuremath{\in} R(L) = RT(W). Applied to W = W\textsubscript{S}: RP\textsuperscript{2} \ensuremath{\to} e \ensuremath{\in} RT(W\textsubscript{S}) \ensuremath{\setminus} W\textsubscript{S}. \ensuremath{\blacksquare}

Remark. Step (c) shows loc(\ensuremath{\Sigma}\textsuperscript{\ensuremath{\infty}} RP\textsuperscript{\ensuremath{\infty}}) = loc(S/2), the 2-torsion spectra; so RT of the reversible localizer generated by RP\textsuperscript{\ensuremath{\infty}} \ensuremath{\to} e is the localizer of Z[1/2]-homology equivalences.


\subsection*{GPT-6 Astra\textquotesingle s reading}

The note supplies the desuspension argument, the localizer-axiom verification and the stable-closure argument. Its decisive outside inputs are the two-local James EHP fibre sequence, Miller's theorem together with the arithmetic-square passage to two-local spheres, Kahn--Priddy in the stated space-retraction form, and Snaith's splitting. Those inputs are named but not proved here; in particular the arithmetic-square passage and the use of iterated bar constructions are compressed. I found no algebraic discrepancy in the displayed cofibre and stable-closure steps, but the precise localization and retraction formulations need independent source checking before this is described as a completed peer review. The source's accessibility discussion is expressly unproved and lies outside the recorded theorem; it is omitted rather than turned into an additional assertion. The relation between $W_S$ and $RT(W_S)$ is exhibited by the specific map $\mathbb{RP}^2\to e$ and the source's explicit construction, not inferred from a categorical label.

\par\bigskip

\section*{8. The two-cone counterexample to the reverse axiom}

\phantomsection\addcontentsline{toc}{section}{8. The two-cone counterexample to the reverse axiom}

\label{nr:entry8}

\noindent\textbf{Statement recorded in the inventory.}\par\smallskip

\textbf{XII Théorème 2 (ms. 78) is false for W\ensuremath{\infty}}: X = A \ensuremath{\star} \{t, s\}, A connected, acyclic, not aspherical. Everything in XII that uses the strong form of W(7 bis) holds for homology localizers only.


\noindent\textbf{Bearing on the manuscript.} XII Th. 4 Cor. 4 and 5, Cor. of ms. 85, Prop. 7, ms. 105, XIIapp ms. 7, 15\par\smallskip

\noindent\textbf{Proved by.} Claude Fable 5.2, effort max (also TOTASPH A.13 of the return)\par

\noindent\textbf{Checked by (inventory record).} Pro return + expert agree\par\smallskip

\noindent\textbf{Cited source.} \path{pro_reviews_max_20261112/DERIVATEURS_EXPERT/FINDINGS.md}, F-G1; \path{pro_reviews_max_20261112/2_TOTAL_ASPHERICITY/REVIEW.md}, A.12--A.13\par\medskip

\subsection*{Statement and proof in the supplied source}

F-G1 (mine). \textbf{XII Théorème 2 (ms. 78), "W\ensuremath{\infty} satisfait à W(7)", is false}, and so is the reverse half of   W(7 bis) for W\ensuremath{\infty}: X = A \ensuremath{\star} \{t, s\} with A connected, acyclic, not aspherical. Grothendieck did not write the   verification ("je devrais sans doute réparer cette omission"). Everything in XII that uses "W(7 bis) forme forte"   (his own footnotes 104, 105 mark these places: Th. 4 Cor. 4 and 5; also the Corollary of ms. 85, Prop. 7, and   the uses at ms. 105, XIIapp ms. 7 and 15) holds for homology localizers and not for W\ensuremath{\infty}. (Also in TOTASPH A.13;   it is the starting point of POLES.)

\subsection*{Proof details recorded in the supplied TOTASPH review}
Own check of the data: 3 + 17 + 15 simplices; arrows 35 + 34 + 45 + 45 = 159; nerve ranks (35, 124, 90). d\textsubscript{2} = [[5, \ensuremath{-}2], [\ensuremath{-}3, 1]], determinant \ensuremath{-}1 (exponent sums of a\textsuperscript{5}b\textsuperscript{-3} and b\textsuperscript{3}(ab)\textsuperscript{-2}). In A\textsubscript{5}: a = (1 4 3 2 5), b = (1 2 3), ab = (1 5)(3 4) (I composed them), so both relators die and a \ensuremath{\ne} 1. Hatcher's Example 2.38 is the binary icosahedral presentation.

X = A \ensuremath{\star} \{u, t\}: U = A \ensuremath{\cup} \{u\}, V = A \ensuremath{\cup} \{t\} are sieves with contractible nerves, BX = \ensuremath{\Sigma}BA is contractible (simply connected and acyclic), U \ensuremath{\cap} V = A is not. I had found the same counterexample independently while reading XII ms. 78. \textbf{This refutes XII Théorème 2 (ms. 78)} as stated; only the direct halves hold. Consequence I add: in II ms. 36ter (Contre-exemple 2.20) the step "i \ensuremath{\in} W\ensuremath{\infty} \ensuremath{\Rightarrow} X/H aspherical" is not available; one gets only that X/H is connected and acyclic. The counterexample 2.20 survives (take H connected with \ensuremath{\Sigma}B(\ensuremath{\Delta}/H) not contractible).


\subsection*{GPT-6 Astra\textquotesingle s reading}

The source's two-cone construction makes the failure concrete: both cones and their union have contractible classifying spaces, while their intersection has the specified nontrivial fundamental group. The supplied TOTASPH review gives the relation matrix and an explicit nontrivial permutation representation for an acyclic example. The suspension is simply connected and acyclic, so its contractibility follows from Hurewicz and Whitehead. This proves the counterexample to the reverse axiom once the category model of the presentation complex is supplied; the review reports that model's counts but the full attaching data are outside the cited FINDINGS paragraph. The list of later manuscript dependencies is a source attribution, not a separate proof of each later statement. The inventory adds the word ``only'' to its homology-localizer sentence; the cited FINDINGS paragraph says ``for homology localizers and not for $W_\infty$.'' The two-cone proof supplies the stated failure for $W_\infty$ but gives no classification excluding every other localizer. The inventory's stronger word is preserved in its exact statement and flagged here as unsupported by that proof.

\par\bigskip

\section*{9. The reported trefoil argument for failure of Loc (7 F)}

\phantomsection\addcontentsline{toc}{section}{9. The reported trefoil argument for failure of Loc (7 F)}

\label{nr:entry9}

\noindent\textbf{Statement recorded in the inventory.}\par\smallskip

\textbf{A direct proof that W\textsubscript{Z} and W\textsubscript{k} fail Loc (7 F)} (trefoil complement over a base of type S\textsuperscript{1}).


\noindent\textbf{Bearing on the manuscript.} XVI ms. 105 (Grothendieck: "Je ne hasarderais aucune conjecture")\par\smallskip

\noindent\textbf{Proved by.} Claude Fable 5.2, effort max\par

\noindent\textbf{Checked by (inventory record).} not yet\par\smallskip

\noindent\textbf{Cited source.} \path{pro_reviews_max_20261112/DERIVATEURS_EXPERT/FINDINGS.md}, F3-5; \path{pro_reviews_max_20261112/3_POLES/REVIEW.md}, C.5\par\medskip

\subsection*{Statement and proof in the supplied source}

\textbf{F3-5 (mine). Direct proof that W\textsubscript{Z} and W\textsubscript{k} fail (7 F)} (trefoil complement over a base of type S\textsuperscript{1};   the Alexander module appears in H\textsubscript{1} of the base change). Independent of Cisinski's properness theory.

\subsection*{The available proof sketch}
\textbf{Independent confirmation for W\textsubscript{Z} and W\textsubscript{k}, k a field} (my argument before reading, \path{work/checks_3/own_thoughts_before_reading.md} item 7, not using E.2): if e \ensuremath{\to} E \ensuremath{\to} S with BS \ensuremath{\simeq} S\textsuperscript{1}, E acyclic and p \ensuremath{\in} Fib\textsubscript{W}, the local system s \ensuremath{\mapsto} H\textsubscript{0}(E/s) must be Z[t, t\textsuperscript{-1}], and base change along a functor Y' \ensuremath{\to} S realizing the map from the trefoil complement to S\textsuperscript{1} gives H\textsubscript{1}(E \ensuremath{\times}\textsubscript{S} Y') \ensuremath{\twoheadrightarrow} Z[t, t\textsuperscript{-1}]/(t\textsuperscript{2} \ensuremath{-} t + 1) \ensuremath{\ne} 0, although E \ensuremath{\times}\textsubscript{S} Y' \ensuremath{\to} E should be a homology isomorphism. This answers XVI 7.5 (ms. 105): \textbf{Loc (6 a, 8, 9) does not imply Loc (7 F)}; and it answers the wish of XVI ms. 124 ("regarder quelques exemples (équivalences homologiques ou cohomologiques, W\textsubscript{n}) \ldots{} si Loc (7 F) est satisfait").


\subsection*{GPT-6 Astra\textquotesingle s reading}

The cited FINDINGS entry gives a claim and one-line description, not a full proof. The supplied POLES review adds the trefoil outline reproduced below, including the Alexander module $\mathbb Z[t,t^{-1}]/(t^2-t+1)$. Its derivation of the fibre local system, the comparison from strict categorical base change to the stated homology calculation, and the claimed surjection are not written out. The review refers for those steps to a working note absent from the archive. Thus the direct proof as supplied has a definite missing middle step. The field-coefficient version is also asserted without its corresponding derivation. This appendix preserves that limitation and does not replace the missing proof with an invented argument.

\par\bigskip

\section*{10. Totality and the fixed-origin path category}

\phantomsection\addcontentsline{toc}{section}{10. Totality and the fixed-origin path category}

\label{nr:entry10}

\noindent\textbf{Statement recorded in the inventory.}\par\smallskip

\textbf{For Y simply connected, Ch(Y; y\textsubscript{0}, \ensuremath{-}) total \ensuremath{\Rightarrow} Y total.}


\noindent\textbf{Bearing on the manuscript.} XII ms. 102--103, 112\par\smallskip

\noindent\textbf{Proved by.} Claude Fable 5.2, effort max\par

\noindent\textbf{Checked by (inventory record).} not yet\par\smallskip

\noindent\textbf{Cited source.} \path{pro_reviews_max_20261112/DERIVATEURS_EXPERT/FINDINGS.md}, F2-11\par\medskip

\subsection*{Statement and proof in the supplied source}

F2-11 (mine). C005 (XII ms. 103) follows from Grothendieck's own remarks: the endpoint functor of   Ch(Y; y\textsubscript{0}, \ensuremath{-}) is final when Y is 1-connected (ms. 102--103), totality passes to final images (Lemma, ms. 112),   and \ensuremath{\Psi} is aspherical and not total (his footnote at ms. 103). General form: for Y simply connected,   Ch(Y; y\textsubscript{0}, \ensuremath{-}) total \ensuremath{\Rightarrow} Y total. Open: the converse, starting with Y = e (the category of interval types).


\subsection*{GPT-6 Astra\textquotesingle s reading}

The proof is the composition of two identified manuscript facts: the endpoint functor is final for the specified simply connected target, and total asphericity passes to final images. The source identifies both manuscript locations, and the supplied review repeats the deduction. These two cited facts establish exactly the displayed implication. Their internal proofs are not reproduced in FINDINGS; this is a short proof by reference rather than a standalone construction of the endpoint category. The converse is not asserted here.

\par\bigskip

\section*{11. Exact colimits, total asphericity, and injective ordinals}

\phantomsection\addcontentsline{toc}{section}{11. Exact colimits, total asphericity, and injective ordinals}

\label{nr:entry11}

\noindent\textbf{Statement recorded in the inventory.}\par\smallskip

\textbf{Grothendieck's conjecture of II ms. 36ter ("totally W\ensuremath{\infty}-aspherical \ensuremath{\Rightarrow} filtered") is false} (\ensuremath{\Delta}\textsubscript{inj}); with the bridge from exactness of abelian colimits to total asphericity (TOTASPH C.5--C.6). The category and the key fact are Isbell--Mitchell (Bull. AMS 79, 1973, Theorem 3); the bridge is not found there.


\noindent\textbf{Bearing on the manuscript.} II ms. 36--36ter\par\smallskip

\noindent\textbf{Proved by.} GPT-6 Astra (effort Pro, ChatGPT Pro mode); proof of C.9 taken from our Open Maths lesson "Derived colimits over a small category" (20 Oct 2026)\par

\noindent\textbf{Checked by (inventory record).} Claude Fable 5.2, effort max\par\smallskip

\noindent\textbf{Cited source.} \path{pro_returns_cited/TOTASPH_sections_C.tex}, C.1--C.10; \path{pro_reviews_max_20261112/2_TOTAL_ASPHERICITY/REVIEW.md}, C.5--C.10\par\medskip

\subsection*{Statement and proof in the supplied source}

\subsection*{An intrinsic criterion and the failure of Lead 4}

\begin{nrdefinition}[Proved: the action criterion]\label{NR:C:action-definition}
Let $R$ be a nonzero commutative ring with identity. We write $\mathcal P_R(A)$ for the following property of a small category $A$: for every covariant functor $H:A\to\Set$, each connected component $E$ of the category of elements $\operatorname{El}_A(H)$ is $R$-acyclic, meaning that $H_0(BE;R)=R$ and $H_j(BE;R)=0$ for $j>0$. No condition is imposed on the empty category of elements. We write $\mathcal P_\infty(A)$ when every such component is weakly contractible. These definitions use covariant, rather than presheaf, categories of elements: an arrow $(a,h)\to(a',h')$ is an arrow $v:a\to a'$ with $H(v)(h)=h'$.
\end{nrdefinition}

\begin{nrlemma}[Proved: homological Theorem A in the final variance]\label{NR:C:homology-A}
Let $f:A\to B$ be a functor. If every $b\backslash A$ is nonempty and $R$-acyclic, then $Bf:BA\to BB$ induces an isomorphism on homology with coefficients in $R$. More generally, if every $b\backslash A$ is nonempty and its reduced $R$-homology vanishes in degrees $0,\ldots,d-1$, for an integer $d\geq1$, then $Bf$ induces isomorphisms in degrees less than $d$ and an epimorphism in degree $d$.
\end{nrlemma}
\begin{proof}
We use the bisimplicial set whose $(p,q)$-simplices are strings
\[
 b_0\longrightarrow\cdots\longrightarrow b_p
 \longrightarrow f(a_0),\qquad
 a_0\longrightarrow\cdots\longrightarrow a_q.
\]
The forgetful map to $NA$ is a weak equivalence after diagonal realization: for a fixed string in $A$, the remaining simplicial direction is the nerve of $B/f(a_0)$, which has terminal object $\id_{f(a_0)}$. The other projection is to $NB$. Taking normalized $R$-chains in both directions and filtering by $p$ gives the spectral sequence of this projection. Its row in degree $q=0$ is the chain complex of $NB$, because each comma category is connected; the rows in degrees $1,\ldots,d-1$ vanish after taking vertical homology. The first possible additional row has degree $d$. Thus the projection induces isomorphisms below degree $d$ and an epimorphism in degree $d$. The first-quadrant filtration is finite in each total degree, so it converges without a boundedness or finiteness assumption on the categories. If all positive rows vanish, it induces isomorphisms in every degree. The two projections commute with the canonical map $Bf$ up to the homotopy supplied by the displayed arrow $b_p\to f(a_0)$; consequently the comparison just computed is the map induced by $f$.
\end{proof}

\begin{nrlemma}[Proved: inheritance by categories of elements]\label{NR:C:inherit-elements}
If $A$ satisfies $\mathcal P_R$, then every connected component $E$ of $\operatorname{El}_A(H)$, for every $H:A\to\Set$, also satisfies $\mathcal P_R$. The corresponding assertion holds for $\mathcal P_\infty$.
\end{nrlemma}
\begin{proof}
The elements lying in $E$ define a subfunctor $H_E\subseteq H$, because arrows preserve connected components. Let $K:E\to\Set$ be any covariant functor. Define
\[
 L(a)=\coprod_{h\in H_E(a)}K(a,h).
\]
For an arrow $v:a\to a'$ in $A$, its restriction to the summand indexed by $h$ is the map $K(v):K(a,h)\to K(a',H(v)h)$, landing in the indicated summand of $L(a')$. This defines a covariant functor $L:A\to\Set$. The map
\[
 ((a,h),k)\longmapsto (a,(h,k))
\]
is an isomorphism of categories $\operatorname{El}_E(K)\cong\operatorname{El}_A(L)$, including their specified arrows. Its connected components therefore have the required property.
\end{proof}

\begin{nrlemma}[Proved: integral acyclicity detects every cyclic subgroup]\label{NR:C:cyclic-covers}
Suppose that a connected, nonempty small category $E$ satisfies $\mathcal P_{\mathbb Z}$. Then $BE$ is weakly contractible. Consequently $\mathcal P_{\mathbb Z}(A)$ and $\mathcal P_\infty(A)$ are equivalent for every small category $A$.
\end{nrlemma}
\begin{proof}
The constant singleton functor on $E$ shows that $BE$ is integrally acyclic. Fix an object $e_0$ and let $\Pi E$ be the groupoid obtained by inverting all arrows of $E$. The usual edge-path presentation of the fundamental group identifies $\operatorname{Aut}_{\Pi E}(e_0)$ with $G=\pi_1(BE,e_0)$. For any subgroup $K\leq G$, define
\[
 H_K(e)=\operatorname{Hom}_{\Pi E}(e_0,e)/K,
\]
where $K$ acts by precomposition, and an arrow $v:e\to e'$ acts by postcomposition with its image in $\Pi E$. All these maps are bijections. The projection $\operatorname{El}_E(H_K)\to E$ is therefore a categorical covering: each arrow has a unique lift with either a prescribed source or a prescribed target. On nerves this gives the usual covering of $BE$ whose monodromy is the transitive coset action of $G$. It is connected, and its fundamental group, at the lift represented by $\id_{e_0}$, is precisely $K$. The covering classification is also recorded in \citet[Theorem~1.38]{Hatcher2002}. These assertions follow directly here by lifting the edge-path presentation; a loop lifts to a loop exactly when its class belongs to $K$.

Take $K$ to be the cyclic subgroup generated by an arbitrary element of $G$. The hypothesis gives
\[
 0=H_1\bigl(B\operatorname{El}_E(H_K);\mathbb Z\bigr)
   =K_{\mathrm{ab}}=K.
\]
The identification of first homology with abelianization is \citet[Theorem~2A.1]{Hatcher2002}. Thus $G$ is trivial. A simply connected acyclic CW complex is weakly contractible: if a first nonzero homotopy group existed, the Hurewicz theorem would identify it with a nonzero homology group \citep[Theorem~4.32]{Hatcher2002}. A nerve realization is a CW complex, so this applies to $BE$.

For general $A$ satisfying $\mathcal P_{\mathbb Z}$, apply Lemma~\ref{NR:C:inherit-elements} and the preceding argument to each component $E$ of each category of elements. The converse follows from homotopy invariance of integral homology.
\end{proof}

\begin{nrtheorem}[Proved: characterization of total asphericity]\label{NR:C:classification}
For every small category $A$, the following conditions are equivalent:
\begin{enumerate}
\item $A$ is totally $\Winfty$-aspherical in the sense of XII, Definition~1.
\item $A$ is totally $W_{\mathbb Z}$-aspherical.
\item $A$ is nonempty and connected, and satisfies $\mathcal P_{\mathbb Z}$.
\item $A$ is nonempty and connected, and satisfies $\mathcal P_\infty$.
\end{enumerate}
In particular, this is a criterion in terms of all covariant set-valued diagrams on $A$, involving no quantification over target categories or final functors. Moreover every connected component of every $\operatorname{El}_A(H)$ is itself totally $\Winfty$-aspherical whenever $A$ is totally $\Winfty$-aspherical.
\end{nrtheorem}
\begin{proof}
The implication (1)$\Rightarrow$(2) follows because a weak homotopy equivalence induces an integral homology isomorphism. Assume (2). The identity functor is final, so $A$ is integrally acyclic and hence nonempty and connected. For any connected nonempty category of elements, the collage of Part~B is a final image of $A$. Its homotopy-pushout description gives the homology exact sequence
\[
 \cdots\to\widetilde H_{j+1}(BA_H;\mathbb Z)
 \to\widetilde H_j(B\operatorname{El}_A H;\mathbb Z)
 \to\widetilde H_j(BA;\mathbb Z)\to\cdots.
\]
Both outside terms vanish. For a component of an arbitrary $H$, first replace $H$ by its component subfunctor. This proves (3). Lemma~\ref{NR:C:cyclic-covers} proves (3)$\Rightarrow$(4).

Assume (4), and let $f:A\to B$ be any ordinary final functor. For $b\in B$, the category $b\backslash A$ is the category of elements of the covariant functor $a\mapsto\operatorname{Hom}_B(b,f(a))$. It is connected and nonempty by finality, and therefore weakly contractible by the hypothesis. Quillen's Theorem~A \citep[Section~1, Theorem~A]{Quillen1973} implies that the specified map $Bf$ is a weak homotopy equivalence. The constant singleton functor in (4) makes $BA$ weakly contractible. Thus $BB$ is weakly contractible, proving (1) by the comprehensive-factorization criterion. The final assertion follows from Lemma~\ref{NR:C:inherit-elements} and (4).
\end{proof}

\begin{nrproposition}[Proved: exact integral colimits are sufficient]\label{NR:C:exact-sufficient}
Let $A$ be nonempty and connected. If $\colim_A:\Ab^A\to\Ab$ is exact, then $A$ is totally $\Winfty$-aspherical.
\end{nrproposition}
\begin{proof}
For $H:A\to\Set$, define $\Sigma_H:\Ab^{\operatorname{El}_A(H)}\to\Ab^A$ by
\[
 (\Sigma_HF)(a)=\bigoplus_{h\in H(a)}F(a,h),
\]
with the transition maps specified in the proof of Lemma~\ref{NR:C:inherit-elements}. It is exact because direct sums of abelian groups are exact. The universal property of a cocone gives the canonical natural isomorphism
\[
 \colim_{\operatorname{El}_A(H)}F
   \cong\colim_A\Sigma_HF.
\]
Thus the colimit functor on each category of elements is exact, and the same is true on any one of its components, by extending diagrams by zero on the other components. The bar resolution computing derived colimits identifies
$L_j\colim_E\underline{\mathbb Z}$ with $H_j(BE;\mathbb Z)$: its term in degree $j$ is the free abelian group on strings of $j$ composable arrows of $E$, and its differential is the alternating sum of the nerve faces. Exactness therefore makes every component $E$ integrally acyclic. Now apply Theorem~\ref{NR:C:classification}.
\end{proof}

\subsection*{An explicit nonfiltered totally aspherical category}

\begin{nrdefinition}[Proved: the category of injective ordinals]\label{NR:C:delta-definition}
Let $\Delta_{\mathrm{inj}}$ have objects $[n]=\{0,\ldots,n\}$ for $n\geq0$, and all strictly increasing functions as arrows, with ordinary composition of functions. This is a countable small category. In the next construction $\mathbb Z[S]$ denotes the free abelian group on a set $S$, and the augmentation of a finite linear combination of arrows is the sum of its coefficients.
\end{nrdefinition}

\noindent\textbf{Historical source wording.} The review corrects the attribution below to Isbell--Mitchell (1973), Theorem~3; the source wording is retained.\par\smallskip
The exact-colimit example is attributed to \citet{Isbell1974Exact} in \citet[Section~2.3, after Theorem~2.5]{Husainov2002Dimension}; its bibliographic record was verified, but the original article's body was not available in this run. We do not use that unavailable body as a proof. The following explicit verification uses polynomial interpolation. A CC0 AI-authored exposition containing the construction was found during the literature search \citep[Section~4]{OpenMaths2026Derived}; we independently verify its algebraic identities here and by the accompanying program \texttt{cd\_affine\_polynomials.py}. Its provenance is recorded in the report, and it is not used as an authority.

\begin{nrlemma}[Proved: interpolation on coordinate hyperplanes]\label{NR:C:interpolation}
Let $F_0,\ldots,F_n$ be polynomials in $\mathbb Z[x_0,\ldots,x_n,u]$ such that $F_i-F_j\in(x_i,x_j)$ for every pair $i,j$. There exists $C$ in this polynomial ring with $C-F_i\in(x_i)$ for every $i$. If all $F_i$ are homogeneous of the same degree, $C$ can be chosen homogeneous of that degree.
\end{nrlemma}
\begin{proof}
For a subset $S\subseteq\{0,\ldots,n\}$ let $\rho_S$ set $x_i=0$ for $i\in S$. For $S\ne\varnothing$, the value $\Phi_S=\rho_S(F_i)$ is independent of the choice of $i\in S$: the difference of two choices belongs to $(x_i,x_j)$ and is killed by $\rho_S$. Set
\[
 C=\sum_{\varnothing\ne S\subseteq\{0,\ldots,n\}}
       (-1)^{|S|+1}\Phi_S.
\]
Modulo $(x_j)$ the terms indexed by $S$ and $S\cup\{j\}$ cancel whenever $j\notin S$ and $S\ne\varnothing$. The remaining term is $\Phi_{\{j\}}$, congruent to $F_j$. All substitutions preserve homogeneity.
\end{proof}

\begin{nrlemma}[Proved: normalized affine cones]\label{NR:C:affine-cones}
For every $n\geq0$ there are an integer $M\geq n$ and
\[
 t\in\mathbb Z[\operatorname{Hom}_{\Delta_{\mathrm{inj}}}([n],[M])]
\]
of augmentation $1$ such that $th=th'$ whenever $h,h':[m]\to[n]$ are parallel arrows. Here composition is extended bilinearly in the free preadditive category. For $n\geq1$ one can take $M=2^{n+1}-2$.
\end{nrlemma}
\begin{proof}
We prove a polynomial assertion simultaneously for all integers $k\geq1$. For variables $z_0,\ldots,z_{n+1}$, let $D_i$ identify $z_i$ and $z_{i+1}$ and relabel the remaining variables in order as $y_0,\ldots,y_n$. We claim that there is a homogeneous integral polynomial $C_{n,k}$ with value $1$ when all variables are $1$, such that $y_i^kD_i(C_{n,k})$ is independent of $i$.

For $n=0$, use $C_{0,k}=1$. Assume the assertion for $n-1$ and all $k$. Let $Q=C_{n-1,2k}$, expressed in the variables $y_0,\ldots,y_n$, and prescribe
\[
 D_i(C_{n,k})=F_i:=Q\prod_{l\ne i}y_l^k.
\]
These prescriptions agree on every pairwise intersection of the hyperplanes $z_i=z_{i+1}$. Indeed, for $i<j$, one must check $D_{j-1}(F_i)=D_i(F_j)$. If $j=i+1$, both omitted-variable products become $\prod_lw_l^k$ and the images of $Q$ agree. If $j\geq i+2$, put $p=i$, $q=j-1$. The two expressions are respectively
\[
 w_q^{2k}D_q(Q)\prod_{l\ne p,q}w_l^k,
 \qquad
 w_p^{2k}D_p(Q)\prod_{l\ne p,q}w_l^k,
\]
and are equal by induction. Choose homogeneous polynomial lifts of the $F_i$. The change of coordinates
$x_i=z_i-z_{i+1}$, $u=z_{n+1}$ is invertible over $\mathbb Z$; the kernels of the $D_i$ are the coordinate ideals $(x_i)$. Lemma~\ref{NR:C:interpolation} therefore supplies $C_{n,k}$. It has value $1$ at the all-ones point, and the required common expression is $Q\prod_l y_l^k$. Its degree satisfies
\[
 d(n,k)=d(n-1,2k)+nk=k(2^{n+1}-n-2).
\]

For an injection $g:[n]\to[M]$, let its gaps be
\[
 e_0=g(0),\quad e_i=g(i)-g(i-1)-1\ (1\leq i\leq n),
 \quad e_{n+1}=M-g(n).
\]
This gives a bijection between injections and monomials of degree $M-n$ in $z_0,\ldots,z_{n+1}$. Extend it linearly, writing $P_t$ for the polynomial of an arrow combination $t$. If $\delta^i:[n-1]\to[n]$ omits $i$, deleting $g(i)$ merges two gaps and the deleted point; hence
\[
 P_{t\delta^i}=y_iD_i(P_t).
\]
Take $P_t=C_{n,1}$ and $M=n+d(n,1)$. Its augmentation is $1$, and all $t\delta^i$ agree.

We finally check that this equalizes every parallel pair, rather than just the cofaces. If their common value is $s$, the relation
$\delta^{b+1}\delta^a=\delta^a\delta^b$ for $a\leq b$ gives $s\delta^a=s\delta^b$. Induction on $n$ therefore shows that $s$ equalizes all parallel injections into $[n-1]$. Every proper injection into $[n]$ factors through a coface; factoring two such injections proves $th=th'$. The identity is the only injection $[n]\to[n]$. The case $n=0$ is immediate.
\end{proof}

\begin{nrtheorem}[Refuted: total asphericity need not imply filteredness]\label{NR:C:counterexample}
The category $\Delta_{\mathrm{inj}}$ is totally $\Winfty$-aspherical and has exact colimits of abelian-group-valued diagrams, but is not filtered. Consequently Lead~4 for $\Winfty$ and manuscript row II-C-003 have a negative answer for arbitrary small categories.
\end{nrtheorem}
\begin{proof}
The affine category $\operatorname{Aff}_{\mathbb Z}(\Delta_{\mathrm{inj}})$, defined and analyzed in Theorem~D.4 of the source return, has the same objects and all integral arrow combinations of augmentation $1$. Any two objects map to a common larger ordinal. Given two parallel affine arrows $\alpha,\beta:[m]\to[n]$, Lemma~\ref{NR:C:affine-cones} provides $t$ for which $th$ is independent of the ordinary injection $h:[m]\to[n]$. Since the coefficients of both $\alpha$ and $\beta$ sum to $1$, we get $t\alpha=t\beta$. Thus the affine category is filtered. Theorem~D.4 of the source return proves exactness of integral colimits. The original category is nonempty and connected, so Proposition~\ref{NR:C:exact-sufficient} proves total $\Winfty$-asphericity.

There are two distinct injections $a,b:[0]\to[1]$, taking $0$ to $0$ and $1$ respectively. Every arrow $c:[1]\to[n]$ is injective, so $ca\ne cb$. Hence this parallel pair has no equalizer by postcomposition, and the category is not filtered. This verifies every hypothesis and the failure in a completely specified countable category; no finite-to-infinite inference is involved.
\end{proof}


\subsection*{The supplied reviewer on attribution and the exact-colimit step}
I verified the return's proof line by line (gap monomials: P\textsubscript{t\ensuremath{\delta}\textsuperscript{i}} = y\textsubscript{i} D\textsubscript{i}(P\textsubscript{t}); the compatibility D\textsubscript{j\ensuremath{-}1}(F\textsubscript{i}) = D\textsubscript{i}(F\textsubscript{j}) in both cases; value 1 at the all-ones point; degree k(2\textsuperscript{n+1} \ensuremath{-} n \ensuremath{-} 2)), and the case n = 1 by hand: t = (1,2) \ensuremath{-} (0,2) + (0,1), t\ensuremath{\delta}\textsuperscript{0} = t\ensuremath{\delta}\textsuperscript{1} = [1]. I also opened the 1973 note (free at ams.org), which the return did not: it states the theorem and sketches a different, combinatorial proof. The label "Proved" should read "Known (Isbell--Mitchell 1973, Th. 3); independent proof".

Aff\textsubscript{Z}(\ensuremath{\Delta}\textsubscript{inj}) filtered \ensuremath{\Rightarrow} colim exact: I checked this direction directly (colim\textsubscript{A} F = colim\textsubscript{Aff A} \ensuremath{\widetilde{F}}, since x \ensuremath{-} (\ensuremath{\Sigma} n\textsubscript{c} c)x = \ensuremath{\Sigma} n\textsubscript{c} (x \ensuremath{-} cx)); it is the "if" part of Isbell--Mitchell's Theorem 1. Then C.6. Not filtered: d\textsuperscript{0}, d\textsuperscript{1}: [0] \ensuremath{\to} [1] are never equalized. \textbf{So the conjecture Grothendieck states at II ms. 36ter is false, and the counterexample is the Isbell--Mitchell category of 1973}; what is new is the identification of total asphericity (C.5) that makes the old example apply. Independent evidence of mine: for several explicit H on \ensuremath{\Delta}\textsubscript{inj} I computed El(H) by hand before reading the proof (subfunctors and quotients of h\textsuperscript{0}, h\textsuperscript{1}; all components came out contractible).


\subsection*{GPT-6 Astra\textquotesingle s reading}

The supplied text writes out the covering-space argument that converts the category-of-elements acyclicity condition into weak contractibility, and it gives the exact direct-sum functor that transfers exact colimits to categories of elements. It then supplies the polynomial affine-cone construction and the two injections which no postcomposition can equalize. The source excerpt retains two dependencies outside its boundaries: the collage model from Part~B and the affine-category exactness theorem from Part~D. The supplied review checks both and gives the short colimit relation used in the latter; its text is reproduced as an accompanying source note. I found the bridge C.5--C.6 and the displayed polynomial identities internally consistent. The theorem that the affine category of injective ordinals is filtered is credited to Isbell--Mitchell, 1973, Theorem~3, as the review requires. This is an independent written proof with recorded dependencies, not a claim of priority for that theorem.

\par\bigskip

\section*{12. The projective-plane counterexample for finite monoids}

\phantomsection\addcontentsline{toc}{section}{12. The projective-plane counterexample for finite monoids}

\label{nr:entry12}

\noindent\textbf{Statement recorded in the inventory.}\par\smallskip

\textbf{Counterexample to Conjecture 7.1 of arXiv:2610.38760v1}: RP\textsuperscript{2} is not BM for a finite monoid M (a finite monoid whose classifying space is finite-dimensional has trivial universal group).


\noindent\textbf{Bearing on the manuscript.} sent to the author on 12 Nov 2026\par\smallskip

\noindent\textbf{Proved by.} GPT-6 Astra (effort Pro)\par

\noindent\textbf{Checked by (inventory record).} Claude Fable 5.2 (effort medium, 11 Nov; again at effort max, 12 Nov)\par\smallskip

\noindent\textbf{Cited source.} \path{pro_returns_cited/CATHOT_sections_G.tex}, G.4 and G.6; \path{pro_reviews_max_20261112/1_CATHOT/G6_VERDICT.md}, theorem and verification\par\medskip

\subsection*{Statement and proof in the supplied source}

\begin{nrsourceclaim}{G.4}{Proved}{A finite monoid cannot have a nonsimply connected finite-dimensional type}
Let $M$ be a finite monoid. Its universal group $U(M)=\pi_1(BM)$ is finite. If $BM$ is homotopy equivalent to a finite-dimensional CW complex, then $U(M)=1$. In particular neither $S^1$ nor $\mathbb{RP}^2$ is the classifying-space type of a finite monoid.
\end{nrsourceclaim}
\begin{proof}
Let $\eta:M\to U(M)$ be the universal homomorphism. For $a\in M$, finiteness gives $a^r=a^s$ for some $r<s$, so $\eta(a)^{s-r}=1$. Every inverse of a generator is therefore a nonnegative power of that generator. It follows that every word in the generators and their inverses equals $\eta(m)$ for some $m\in M$. Thus $\eta$ is surjective, and $U(M)$ is finite. The identification with $\pi_1(BM)$ follows directly from the edge-and-triangle presentation: the generators are the nonidentity elements and the relations are $[a][b]=[ab]$.

Choose an idempotent $e$ in the minimal nonempty two-sided ideal $I$ of $M$. Such an ideal exists by finiteness, and a sufficiently large power of any of its elements is idempotent. We verify that $H=eMe$, with identity $e$, is a group. For $x\in eMe$, minimality gives $MxM=I$, so $e=axb$ for some $a,b\in M$. Multiplying by $e$ at both ends yields
\[
                    e=(eae)x(ebe)
\]
in the finite monoid $eMe$. In any finite monoid a product is invertible only if each factor is invertible: a right inverse makes right multiplication injective, hence surjective, and gives a two-sided inverse; the dual argument treats a left inverse. Thus $x$ is invertible. The group $H$ is finite, and the homomorphism $H\to U(M)$ is surjective since $\eta(eme)=\eta(m)$.

There is a canonical morphism $BH\to BM$ in $\Hot$ even though the inclusion $H\subset M$ sends $1_H$ to $e$, not necessarily to $1_M$. To type this map, let $\operatorname{Kar}(M)$ be the category of idempotents of $M$ with
\[
       \Hom_{\operatorname{Kar}(M)}(f,g)=gMf.
\]
The full one-object subcategory at $1_M$ is $M$ and the full one-object subcategory at $e$ is $H$. Every object is a retract of $1_M$. The inclusion $j:M\to\operatorname{Kar}(M)$ is a weak equivalence: for each object $f$, the comma category $j/f$ is a retract of $j/1_M$, which has a terminal object. A retract of a contractible space is contractible, and Quillen's Theorem A applies. Hence the typed morphism is
\[
          BH\longrightarrow B\operatorname{Kar}(M)
                     \xleftarrow[\simeq]{Bj} BM.
\]
It induces the surjection $H\to U(M)$ on fundamental groups, after the basepoint path given by the retract arrows is chosen.

Suppose $h:BM\to Z$ is a homotopy equivalence with $Z$ a connected finite-dimensional CW complex. Miller's theorem states that every map from the classifying space of a torsion group to $Z$ induces the zero map on fundamental groups \citep[Theorem B, p.~40]{Miller1984}. Apply it to the composite $BH\to BM\xrightarrow{h}Z$. Its map on fundamental groups is both zero and surjective. Consequently $\pi_1(Z)=U(M)=1$. Notice that this step uses the precise fundamental-group theorem, with no assumption that $Z$ was simply connected.
\end{proof}

\begin{nrsourceclaim}{G.6}{Refuted}{Finite fundamental group is not sufficient for finite monoids}
The assertion that every connected finite CW complex with finite fundamental group is the classifying space of a finite monoid is false.
\end{nrsourceclaim}
\begin{proof}
Take $\mathbb{RP}^2$, whose fundamental group is $\Z/2$. G.4 excludes any finite monoid realization. In particular this refutes Conjecture 7.1 in the verified version of \citet{BenBassat2026}; it does not refute that paper's simply connected Theorem C. The explicit finite categorical and poset models below verify that the obstruction concerns the one-object restriction.
\end{proof}
\subsection*{The supplied verdict	extquotesingle s verification of the finite-group step}
\textbf{The proof of step 4, so that nothing is taken on authority.} Let f: BG \ensuremath{\to} X with G finite, X a connected CW complex of dimension d, and suppose f\textsubscript{*} = \ensuremath{\varphi} is non-trivial on \ensuremath{\pi}\textsubscript{1}. Pick c \ensuremath{\in} G of prime-power order p\textsuperscript{k} with \ensuremath{\varphi}(c) \ensuremath{\ne} 1 (a primary component of an element not in the kernel), C = \ensuremath{\langle}c\ensuremath{\rangle}, N = \ensuremath{\varphi}(C) \ensuremath{\ne} 1, X\ensuremath{\prime} \ensuremath{\to} X the covering with \ensuremath{\pi}\textsubscript{1}(X\ensuremath{\prime}) = N, f\ensuremath{\prime}: BC \ensuremath{\to} X\ensuremath{\prime} the lift, \ensuremath{\psi} a character of N of order p, L the flat line bundle of \ensuremath{\psi} on X\ensuremath{\prime}. Then f\ensuremath{\prime}\textsuperscript{*}L is the line bundle of \ensuremath{\chi} = \ensuremath{\psi}\ensuremath{\circ}\ensuremath{\varphi}|C on BC, and \ensuremath{\chi} has order p. In K\textsuperscript{0}(X\ensuremath{\prime}), ([L] \ensuremath{-} 1)\textsuperscript{d+1} = 0, because the skeletal filtration is multiplicative and X\ensuremath{\prime} has dimension d (Miller argues with Lusternik--Schnirelmann category). So (\ensuremath{\chi} \ensuremath{-} 1)\textsuperscript{d+1} restricts to zero in K\textsuperscript{0} of every finite skeleton of BC. By Atiyah, \emph{Characters and cohomology of finite groups}, Publ. Math. IHES 9 (1961), Proposition (8.1) for cyclic groups (or Theorem (7.2), p. 46: \ensuremath{\widehat{\alpha}}: R(G)\textasciicircum{} \ensuremath{\to} \ensuremath{\mathcal K}*(B\textsubscript{G}) is a topological isomorphism) and Proposition (6.11), p. 45 (for a p-group R(G) \ensuremath{\to} R(G)\textasciicircum{} is a monomorphism), (\ensuremath{\chi} \ensuremath{-} 1)\textsuperscript{d+1} = 0 in R(C). But its value at c is (\ensuremath{\zeta} \ensuremath{-} 1)\textsuperscript{d+1} \ensuremath{\ne} 0 with \ensuremath{\zeta} a primitive p-th root of unity. Contradiction. (Miller's own computation: K\textsuperscript{0}(BZ\textsubscript{mn}) = Z[[x]]/((1+x)\textsuperscript{mn} \ensuremath{-} 1) and the image (1+x)\textsuperscript{n} \ensuremath{-} 1 of x\textsubscript{m} is not nilpotent; I verified this by mapping x to \ensuremath{\zeta} \ensuremath{-} 1 in Z\textsubscript{p}[\ensuremath{\zeta}] for \ensuremath{\zeta} of order p\textsuperscript{c}, p\textsuperscript{c} the exact power of a prime p | m in mn.)


\subsection*{GPT-6 Astra\textquotesingle s reading}

The proof specifies the exact morphism from the maximal subgroup's classifying space to that of the monoid through its Karoubi envelope, and computes the induced surjection on fundamental groups. Miller's stated theorem makes the same induced map trivial for a finite-dimensional target; the two conclusions force $U(M)=1$. The fundamental group of $\mathbb{RP}^2$ then supplies the counterexample. I found the source proof complete from Quillen's Theorem~A and Miller's stated theorem. The supplied verdict adds a $K$-theoretic verification of the latter finite-group case. The separately cited EMAIL\_CLAIMS\_CHECK file is absent, so its purported additional checks cannot be reproduced here. The refutation is tied to the conjecture as quoted in the supplied version-1 record; no assertion about a later version or publication status is added.

\par\bigskip


\section*{Statements found to be known}
\phantomsection\addcontentsline{toc}{section}{Statements found to be known}
The inventory closes by recording the following prior references, so that these statements are not presented as priority claims.
\begin{enumerate}
\item The filteredness of the affine category of injective ordinals: J.~R.~Isbell and B.~Mitchell, ``Exact colimits,'' \emph{Bulletin of the American Mathematical Society} \textbf{79} (1973), 994--996, Theorem~3. This is the prior theorem used in section~11; the bridge to total asphericity is the additional connection recorded there. The source review also identifies the intermediate CC0 Open Maths lesson, ``Derived colimits over a small category'' (20 October 2026), as the provenance of the return's polynomial proof.
\item Uniqueness in Nielsen realization: W.~J.~Harvey, ``On branch loci in Teichmüller space,'' \emph{Transactions of the American Mathematical Society} \textbf{153} (1971), 387--399. NEW\_RESULTS records that the La Longue Marche expert located this reference; its proof and precise formulation are outside the supplied new-results proof files.
\item The identification recorded as ``localiseur spécial'' and accessible fundamental localizer: D.-C.~Cisinski, \emph{Les préfaisceaux comme modèles des types d'homotopie}, Astérisque~308 (2006), Corollary~4.2.22, as cited in NEW\_RESULTS and the supplied expert FINDINGS F3-4. The exact axiom translation is not added as a new appendix theorem.
\end{enumerate}
\endgroup

